% Options for packages loaded elsewhere
\PassOptionsToPackage{unicode}{hyperref}
\PassOptionsToPackage{hyphens}{url}
\documentclass[
  ukrainian,
]{article}
\usepackage{xcolor}
\usepackage[margin=2.5cm]{geometry}
\usepackage{amsmath,amssymb}
\setcounter{secnumdepth}{-\maxdimen} % remove section numbering
\usepackage{iftex}
\ifPDFTeX
  \usepackage[T1]{fontenc}
  \usepackage[utf8]{inputenc}
  \usepackage{textcomp} % provide euro and other symbols
\else % if luatex or xetex
  \usepackage{unicode-math} % this also loads fontspec
  \defaultfontfeatures{Scale=MatchLowercase}
  \defaultfontfeatures[\rmfamily]{Ligatures=TeX,Scale=1}
\fi
\usepackage{lmodern}
\ifPDFTeX\else
\setmainfont{DejaVu Serif}
\setsansfont{DejaVu Sans}
\setmonofont{DejaVu Sans Mono}
\setmathfont{STIX Two Math}
\fi
\ifPDFTeX\else
  % xetex/luatex font selection
\fi
% Use upquote if available, for straight quotes in verbatim environments
\IfFileExists{upquote.sty}{\usepackage{upquote}}{}
\IfFileExists{microtype.sty}{% use microtype if available
  \usepackage[]{microtype}
  \UseMicrotypeSet[protrusion]{basicmath} % disable protrusion for tt fonts
}{}
\makeatletter
\@ifundefined{KOMAClassName}{% if non-KOMA class
  \IfFileExists{parskip.sty}{%
    \usepackage{parskip}
  }{% else
    \setlength{\parindent}{0pt}
    \setlength{\parskip}{6pt plus 2pt minus 1pt}}
}{% if KOMA class
  \KOMAoptions{parskip=half}}
\makeatother
\usepackage{graphicx}
\graphicspath{{/Users/kalinicenkovitalijmikolajovic/.codex/worktrees/5c15/Script_for_translate_book/Output_Final/}}
\setkeys{Gin}{width=0.82\linewidth,height=0.82\textheight,keepaspectratio}
\makeatletter
\newsavebox\pandoc@box
\newcommand*\pandocbounded[1]{% scales image to fit in text height/width
  \sbox\pandoc@box{#1}%
  \Gscale@div\@tempa{\textheight}{\dimexpr\ht\pandoc@box+\dp\pandoc@box\relax}%
  \Gscale@div\@tempb{\linewidth}{\wd\pandoc@box}%
  \ifdim\@tempb\p@<\@tempa\p@\let\@tempa\@tempb\fi% select the smaller of both
  \ifdim\@tempa\p@<\p@\scalebox{\@tempa}{\usebox\pandoc@box}%
  \else\usebox{\pandoc@box}%
  \fi%
}
% Set default figure placement to htbp
\def\fps@figure{htbp}
\makeatother
\ifLuaTeX
\usepackage[bidi=basic,shorthands=off,provide=*]{babel}
\else
\usepackage[bidi=default,shorthands=off,provide=*]{babel}
\fi
\ifLuaTeX
  \usepackage{selnolig} % disable illegal ligatures
\fi
\setlength{\emergencystretch}{3em} % prevent overfull lines
\providecommand{\tightlist}{%
  \setlength{\itemsep}{0pt}\setlength{\parskip}{0pt}}
\usepackage{bookmark}
\IfFileExists{xurl.sty}{\usepackage{xurl}}{} % add URL line breaks if available
\urlstyle{same}
\hypersetup{
  pdflang={uk},
  hidelinks,
  pdfcreator={LaTeX via pandoc}}

\author{}
\date{}

\begin{document}

ЧАСТКОВО ДИФЕРЕНЦІАЛ РІВНЯННЯ для Вчені та Інженери

Стенлі Дж. Фарлоу, професор математики, Університет Мену

\section{С. Фарлоу}\label{ux441.-ux444ux430ux440ux43bux43eux443}

\section{РІВНЯННЯ З ПОХІДНИМИ ПОХІДНИМИ ДЛЯ НАУКОВЦІВ І
ІНЖЕНЕРІВ}\label{ux440ux456ux432ux43dux44fux43dux43dux44f-ux437-ux43fux43eux445ux456ux434ux43dux438ux43cux438-ux43fux43eux445ux456ux434ux43dux438ux43cux438-ux434ux43bux44f-ux43dux430ux443ux43aux43eux432ux446ux456ux432-ux456-ux456ux43dux436ux435ux43dux435ux440ux456ux432}

Перекладено з англійської \textbf{A.I.} Плісом

Редактор

С. І. Похожаєва

5\textbar{} 66K 22.161.6 φ24 υ ΥΠΚ 517.2

\section{Фарлоу С.}\label{ux444ux430ux440ux43bux43eux443-ux441.}

F 24 Рівняння з частими похідними для науковців та інженерів. Москва,
видавництво Мир, 1985, 384 с.

King's American Mathematician, який є підручником з теорії
диференціальних рівнянь у похідних похідних. Він характеризується
компактністю, ясністю та ясністю подачі, а також неформальним підходом
до подачі матеріалу. Вона містить багато ілюстрацій, графіків і діаграм,
але замість суворих доказів часто даються роздуми, засновані на інтуїції
або аналогіях.

Для інженерів і нематематиків --- біологів, диміків, а також студентів
університету.
\texttt{\textbackslash{}Phi\ =\ \textbackslash{}frac\{1702050000-380\}\{041(01)-85\}\ 26-85,\ \textbackslash{}text\{\ q.\ 1\}}
BBK 22.161.6

Редакційна колегія журналу Literature on Mathematical Sciences

© переклад російською, Мир, 1985

© John Wiley \& Sons, Inc.~1982. Всі права захищені. Авторизований
переклад з англомовного видання, опублікований John Wiley \& Sons, Inc.

\section{З РЕДАКТОРА
ПЕРЕКЛАДУ}\label{ux437-ux440ux435ux434ux430ux43aux442ux43eux440ux430-ux43fux435ux440ux435ux43aux43bux430ux434ux443}

Книга, запропонована для читачів, слід розглядати як підручник для
курсів: «Диференціальні рівняння в похідних похідних» та «Рівняння
математичної фізики» для студентів технічних університетів. Вона також
буде корисною для інженерів різних спеціальностей. У книзі представлено
традиційні методи розв'язання диференціальних рівнянь у похідних
похідних, включаючи деякі чисельні методи. Водночас особлива увага
приділяється методу розділення змінних і методу інтегральних
перетворень, і цю частину книги можна використовувати на практичних
заняттях в університетах

Орієнтація книги на широку «нематематичну» аудиторію визначила стиль
викладання розглянутих питань: інтуїтивний підхід і велика увага до
фізичного значення не лише самих рівнянь, а й меж і початкових умов для
різних задач. Вважається, що читачі знайомі лише з основами
диференціального та інтегрального числення та елементами теорії
звичайних диференціальних рівнянь.

Фахівці можуть не задовольнитися інтуїтивним описом певних понять і
методів автором, але більш суворий виклад неминуче призведе до
ускладнення книги і звузить коло її читачів, що, очевидно, не було
наміром автора.

Більшим запереченням може бути відсутність загальних результатів щодо
задач для диференціальних рівнянь у похідних з змінними коефіцієнтами,
включно з рівняннями вищого порядку. Однак читач має чітко розуміти, що
ця книга є вступним курсом з теорії диференціальних рівнянь у часткових
похідних. Питання лінійної теорії диференціальних рівнянь, які автор не
включив, можна знайти в літературі, наведеній наприкінці книги.

Матеріал книги подається автором з великою методологічною майстерністю.
Однак не все було однаково успішним, і винний читач знайде це,
наприклад, у вивченні принципу Дюамеля, сферичних заповітів та деяких
інших уривків у книзі.

Без сумніву, запропонована книга С. Фарлоу знайде читача не лише серед
студентів і інженерів, а й серед викладачів, які можуть багато чого
навчитися з конкретного способу презентації та розробки спеціальних
математичних курсів.

\section{ПЕРЕДМОВА}\label{ux43fux435ux440ux435ux434ux43cux43eux432ux430}

За останні кілька років спостерігається значне зростання кількості
учнів, які починають вивчати диференціальні рівняння з похідними у своїх
одинадцятих класах. Крім того, багато з них спеціалізуються на сферах,
де інтуїція важливіша за математичну строгість. Тому, пишучи цю книгу, я
прагнув створити текст, який стимулює інтуїтивне мислення читачів, не
втрачаючи при цьому значної математичної строгості. Можливо, на одному
полюсі, подати предмет на високому рівні «епсилон-дельта», але тоді, як
це зазвичай буває, студенти не знатимуть, що робити з отриманою
інформацією. Інша крайність --- повне ігнорування математичних тонкощів,
але тоді ні учні, ні вчителі не знатимуть, як діяти. Я намагався знайти
розумний компроміс між цими крайнощами.

Ця книга виникла з курсу лекцій, які я читав протягом останніх п'яти
років. Розташування матеріалу не є загальноприйнятим --- вони майже
незалежні один від одного

Друг лекцій.

Основна увага приділяється двом найважливішим аналітичним методам:
розділенню змінних та інтегральним перетворенням. Також представлені
деякі нестандартні розділи, такі як метод Монте-Карло, варіаційне
числення, теорія керування, теорія потенціалів, інтегральні рівняння. Це
пов'язано з тим, що більшості студентів доведеться стикатися з цими
питаннями під час навчання, і якщо вони не ознайомляться з ними зараз,
навряд чи коли-небудь систематично їх вивчатимуть.

Цю книгу можна використовувати як навчальний посібник для одно- або
двосеместрового курсу як для початківців, так і для старших студентів.
Читачі повинні знати лише диференціальне та інтегральне числення та
звичайні диференціальні рівняння. Матеріал більшості лекцій подається в
одному-двох уроках. Типовий односеместровий курс може включати лекції
1-13, 15-17, 19-20, 22-23, 25-27, 30-32, 37-39. Усі 47 лекцій легко
читати протягом двох семестрів. Водночас часу буде достатньо для
вирішення проблем.

Автор вдячний Вайлі за пропозицію написати книгу, а також рецензентам
--- професорам Крісу Рорресу та М. Куршиду Алі, які зробили низку
корисних коментарів до тексту книги. Будь-які пропозиції щодо подальшого
покращення цього підручника будуть вдячні як від учнів, так і від
викладачів. Я також вдячний Дороті, Сюзан, Александер і Дезі Фарлоу.

Стенлі Дж. Фарлоу

\section{ВСТУП ДО ТЕОРІЇ ДИФЕРЕНЦІАЛЬНИХ РІВНЯНЬ У ПОХІДНИХ
ПОХІДНИХ}\label{ux432ux441ux442ux443ux43f-ux434ux43e-ux442ux435ux43eux440ux456ux457-ux434ux438ux444ux435ux440ux435ux43dux446ux456ux430ux43bux44cux43dux438ux445-ux440ux456ux432ux43dux44fux43dux44c-ux443-ux43fux43eux445ux456ux434ux43dux438ux445-ux43fux43eux445ux456ux434ux43dux438ux445}

МЕТА ЛЕКЦІЇ: Показати, що таке диференціальні рівняння в похідних
похідних, де і як вони виникають і як їх розв'язують. Коротко розгляньте
методи класифікації рівнянь і надайте список базових понять, які будуть
детально вивчені на наступних лекціях.

Більшість фізичних явищ у таких сферах, як гідродинаміка, електрика та
магнетизм, механіка, оптика та теплопередача, загалом можна описати за
допомогою рівнянь у похідних похідних (PNP). Більшість рівнянь у
математичній фізиці є диференціальними рівняннями з похідними похідними.
Правда, з деякими спрощеними припущеннями ці рівняння зводяться до
звичайних диференціальних рівнянь, але повний опис таких систем неминуче
призводить до використання диференціальних рівнянь у похідних похідних.

\section{Що таке часткові диференціальні
рівняння?}\label{ux449ux43e-ux442ux430ux43aux435-ux447ux430ux441ux442ux43aux43eux432ux456-ux434ux438ux444ux435ux440ux435ux43dux446ux456ux430ux43bux44cux43dux456-ux440ux456ux432ux43dux44fux43dux43dux44f}

Рівняння з частковими похідними --- це рівняння, яке містить
диференціали у похідних похідних. На відміну від звичайних
диференціальних рівнянь (TAC), де невідома функція залежить лише від
однієї змінної, у рівняннях часткових похідних вона залежить від кількох
змінних (наприклад, температура \texttt{u(x,\ t)} залежить від
координати x і часу t).

Тепер наведемо деякі з найважливіших диференціальних рівнянь у похідних
похідних. Для спрощення запису ми використаємо наступні нотації

\begin{verbatim}
u_t \equiv \frac{\partial u}{\partial t}
\end{verbatim}

\texttt{,\ \textquotesingle{}u\_x\ \textbackslash{}equiv\ \textbackslash{}frac\{\textbackslash{}partial\ u\}\{\textbackslash{}partial\ x\}\textquotesingle{}\ ,\ \textquotesingle{}u\_\{xx\}\ \textbackslash{}equiv\ \textbackslash{}frac\{\textbackslash{}partial\^{}2\ u\}\{\textbackslash{}partial\ x\^{}2\}\textquotesingle{}\ ,\ ...}

\section{Деякі диференціальні рівняння в похідних
похідних}\label{ux434ux435ux44fux43aux456-ux434ux438ux444ux435ux440ux435ux43dux446ux456ux430ux43bux44cux43dux456-ux440ux456ux432ux43dux44fux43dux43dux44f-ux432-ux43fux43eux445ux456ux434ux43dux438ux445-ux43fux43eux445ux456ux434ux43dux438ux445}

\begin{verbatim}
u_t = u_{xx}
\end{verbatim}

(одновимірне рівняння теплопровідності),
\texttt{u\_t\ =\ u\_\{xx\}\ +\ u\_\{yy\}} (двовимірне рівняння
теплопровідності),
\texttt{u\_\{rr\}\ +\ \textbackslash{}frac\{1\}\{r\}u\_r\ +\ \textbackslash{}frac\{1\}\{r\^{}2\}u\_\{00\}\ =\ 0}
(рівняння Лапласа за полярними координатами),
\texttt{u\_\{tt\}\ =\ u\_\{xx\}\ +\ u\_\{yy\}\ +\ u\_\{zz\}} (тривимірне
хвильове рівняння),
\texttt{u\_\{tt\}\ =\ u\_\{xx\}\ +\ zu\_t\ +\ \textbackslash{}beta\ u}
(телеграфне рівняння).

\section{Примітки щодо наведених
прикладів}\label{ux43fux440ux438ux43cux456ux442ux43aux438-ux449ux43eux434ux43e-ux43dux430ux432ux435ux434ux435ux43dux438ux445-ux43fux440ux438ux43aux43bux430ux434ux456ux432}

У всіх наведених вище прикладах функція невідома і залежить від більше
ніж однієї змінної. Змінна і (яку ми диференціюємо) називається залежною
змінною. Змінні, за якими відбувається диференціювання, називаються
незалежними змінними. Наприклад, у рівнянні.

\begin{verbatim}
u_t = u_{xx}
\end{verbatim}

Залежна змінна \texttt{u(x,\ t)} є функцією двох незалежних змінних x і
t; та в рівнянні

\begin{verbatim}
u_t = u_{rr} + \frac{1}{r} u_r + \frac{1}{r^2} u_{00}
\end{verbatim}

змінна \texttt{u(r,\ \textbackslash{}theta,\ t)} залежить від
\texttt{r,\ \textbackslash{}theta,\ t} .

\section{Чому необхідно вивчати диференціальні рівняння у похідних
похідних?}\label{ux447ux43eux43cux443-ux43dux435ux43eux431ux445ux456ux434ux43dux43e-ux432ux438ux432ux447ux430ux442ux438-ux434ux438ux444ux435ux440ux435ux43dux446ux456ux430ux43bux44cux43dux456-ux440ux456ux432ux43dux44fux43dux43dux44f-ux443-ux43fux43eux445ux456ux434ux43dux438ux445-ux43fux43eux445ux456ux434ux43dux438ux445}

Більшість фізичних законів природи можна сформулювати через рівняння з
похідними похідними. Прикладами є рівняння Максвелла, закон
теплопередачі Ньютона, рівняння Нав'є-Стокса, рівняння руху Ньютона та
рівняння Шредінгера в кватичній механіці. У всіх цих рівняннях фізичні
явища описуються мовою просторових і часових похідних. Похідні
з'являються в рівняннях, оскільки вони описують найважливіші фізичні
величини (такі як швидкість, прискорення, сила, тертя, потік, струм
тощо). Таким чином, виникають диференціальні рівняння в частинних
похідних, які містять невідому функцію, яку потрібно визначити.

Мета цієї книги --- показати читачеві:

\begin{enumerate}
\def\labelenumi{\arabic{enumi}.}
\item
  Як сформулювати фізичну задачу у формі диференціального рівняння у
  похідних похідних (побудова математичної моделі).
\item
  Як розв'язати диференціальне рівняння в похідних (з урахуванням
  початкових і граничних умов).
\end{enumerate}

Перш ніж перейти до побудови математичних моделей, давайте коротко
розглянемо методи розв'язання рівнянь у похідних похідних.

\section{Як розв'язати диференціальні рівняння в похідних
похідних?}\label{ux44fux43a-ux440ux43eux437ux432ux44fux437ux430ux442ux438-ux434ux438ux444ux435ux440ux435ux43dux446ux456ux430ux43bux44cux43dux456-ux440ux456ux432ux43dux44fux43dux43dux44f-ux432-ux43fux43eux445ux456ux434ux43dux438ux445-ux43fux43eux445ux456ux434ux43dux438ux445}

Гарне питання. Виявляється, існує цілий арсенал методів, придатних для
практичного використання. Найважливішими є ті, у яких диференціальні
рівняння з похідними похідними зводяться до звичайних диференціальних
рівнянь. Перелічимо десять методів розв'язання диференціальних рівнянь у
похідних похідних:

\begin{enumerate}
\def\labelenumi{\arabic{enumi}.}
\tightlist
\item
  Метод розділення змінних. Рівняння з факторами, отриманими з n
  незалежними змінними, зводиться до n звичайних-
\end{enumerate}

диференціальнірівняння 1).

\begin{enumerate}
\def\labelenumi{\arabic{enumi}.}
\setcounter{enumi}{1}
\item
  Метод інтегральних перетворень. Рівняння в похідних з n незалежними
  змінними зводиться до рівняння з похідними у (n-1) незалежних змінних;
  Відповідно, рівняння з частковими похідними з двома незалежними
  змінними можна звести до звичайного диференціального рівняння
  \texttt{\^{}2} ).
\item
  Метод перетворення координат. Початкове диференціальне рівняння в
  похідних похідних зводиться до звичайного рівняння з частковими
  похідними або до іншого, простішого рівняння у похідних похідних
  шляхом відповідного перетворення координат (наприклад, обертання осей
  координат тощо).
\item
  Перетворення залежної змінної. Початкове диференціальне рівняння в
  похідних похідних перетворюється у таке рівняння з похідними похідними
  для іншої невідомої функції, яка
\end{enumerate}

легше розв'язати, ніж початкові 3).

\begin{enumerate}
\def\labelenumi{\arabic{enumi}.}
\setcounter{enumi}{4}
\tightlist
\item
  Чисельні методи. Початкове диференціальне рівняння в похідних похідних
  зводиться до системи різницей, яку розв'язують методом ітерацій на
  комп'ютері. У багатьох випадках це єдиний спосіб розв'язати
  диференціальне рівняння з похідними похідними. Окрім різницьких
  методів для розв'язання рівнянь з частковими про-
\end{enumerate}

\begin{enumerate}
\def\labelenumi{\alph{enumi})}
\setcounter{enumi}{4}
\tightlist
\item
  Ефективніше використовувати загальні перетворення незалежних і
  залежних змінних, включаючи похідні. Ед.
\end{enumerate}

4) Це стосується повного розділення змінних. При частковому розділенні
змінних одне диференціальне рівняння в похідних похідних зводиться до
кількох рівнянь у часткових похідних з меншою кількістю незалежних
змінних. --- Ед. Ед.

2) Це відбувається у випадку одномірного інтегрального перетворення. У
випадку k-вимірного інтегрального перетворення LCP з n незалежними
змінними зводиться до рівняння часткових з n---k незалежними змінними.
--- Ед. Ед.

Існують також інші чисельні методи, зокрема ті, що базуються на
наближенні розв'язків за поліноміальними поверхнями (наближення за
сплайнами).

\begin{enumerate}
\def\labelenumi{\arabic{enumi}.}
\setcounter{enumi}{5}
\tightlist
\item
  Метод теорії збурень. Початкова нелінійна задача зводиться до
  послідовності лінійних задач, апрокенн-
\end{enumerate}

нелінійна задача 1).

\begin{itemize}
\tightlist
\item
  \begin{enumerate}
  \def\labelenumi{\arabic{enumi}.}
  \setcounter{enumi}{6}
  \tightlist
  \item
    Метод функцій Гріна. Початкові та граничні умови замінюються
    системою найпростіших джерел, і задача розв'язується для кожного
    найпростішого джерела. Повне розв'язання початкової задачі
    отримується шляхом підсумовування розв'язків елементарних джерел.
  \end{enumerate}
\item
  \begin{enumerate}
  \def\labelenumi{\arabic{enumi}.}
  \setcounter{enumi}{7}
  \tightlist
  \item
    Метод інтегральних рівнянь. Рівняння з частковими похідними
    зводяться до інтегрального рівняння (рівняння, в якому невідома
    функція знаходиться під знаком інтеграла). Існує багато різних
    методів розв'язання інтегральних рівнянь.
  \end{enumerate}
\item
  \begin{enumerate}
  \def\labelenumi{\arabic{enumi}.}
  \setcounter{enumi}{8}
  \tightlist
  \item
    Варіаційні методи. Замість рівняння з частковими похідними
    диференціальними розв'язують задачу мінімізації. Виявляється, що
    функція, яка дає мінімум певному виразу (наприклад, загальній
    енергії системи), одночасно є розв'язком початкового рівняння
    похідниху похідних 2.
  \end{enumerate}
\item
  \begin{enumerate}
  \def\labelenumi{\arabic{enumi}.}
  \setcounter{enumi}{9}
  \tightlist
  \item
    Метод розкладу за власними функціями. Розв'язок рівняння часткових
    похідних шукається як послідовність власних функцій. Ці власні
    функції знаходяться як розв'язки так званої задачі власних значень,
    що відповідає початковій задачі для рівняння часткових похідних.
  \end{enumerate}
\end{itemize}

\section{Типи рівнянь у похідних
похідних}\label{ux442ux438ux43fux438-ux440ux456ux432ux43dux44fux43dux44c-ux443-ux43fux43eux445ux456ux434ux43dux438ux445-ux43fux43eux445ux456ux434ux43dux438ux445}

Рівняння в похідних похідних можна класифікувати за багатьма
характеристиками. Класифікація рівнянь важлива, оскільки, як виявилося,
кожен клас має свою загальну теорію та методи розв'язання рівнянь.

Тут ми наведемо шість основних методів класифікації

рівняння.

\begin{enumerate}
\def\labelenumi{\arabic{enumi}.}
\tightlist
\item
  Порядок рівняння. Порядок рівняння --- це найвищий порядок часткових
  похідних, включених до рівняння. Наприклад,
\end{enumerate}

Примітка. Ед.

1) Точніше, це метод лінеаризації нелінійних задач. Загальна теорія
збурень містить лінійні та нелінійні теорії збурень, кожна з яких, у
свою чергу, є регулярною та особливою. --- Ед. Ед.

Ці методи застосовні до спеціальних задач, що зумовлено існуванням
відповідних виразів --- функціоналів для цих задач.

Застосовувався раніше

\begin{verbatim}
u_t = u_{xx}
\end{verbatim}

(рівняння другого порядку), \texttt{u\_t\ =\ u\_x} (рівняння першого
порядку), \texttt{u\_1\ =\ uu\_\{xxx\}\ +\ \textbackslash{}sin\ x}
(рівняння трьох порядків).

\begin{enumerate}
\def\labelenumi{\arabic{enumi}.}
\setcounter{enumi}{1}
\tightlist
\item
  Кількість змінних. Кількість змінних --- це кількість незалежних
  змінних. Наприклад,
\end{enumerate}

\begin{verbatim}
u_t = u_{xx}
\end{verbatim}

\texttt{(рівняння\ з\ двома\ змінними\ \textquotesingle{}x\textquotesingle{}\ і\ \textquotesingle{}t\textquotesingle{}\ ),\ \textquotesingle{}u\_t\ =\ u\_\{rr\}\ +\ \textbackslash{}frac\{1\}\{r\}\ u\_r\ +\ \textbackslash{}frac\{1\}\{r\^{}2\}\ u\_\{\textbackslash{}theta\textbackslash{}theta\}\textquotesingle{}\ (рівняння\ з\ трьома\ змінними\ \textquotesingle{}r\textquotesingle{}\ ,}

\begin{enumerate}
\def\labelenumi{\arabic{enumi}.}
\setcounter{enumi}{2}
\tightlist
\item
  Лінійність. Рівняння з частковими похідними можуть бути лінійними та
  нелінійними. У лінійних рівняннях залежна змінна та всі її часткові
  похідні включаються лінійно, зокрема, вони не множаться одна на одну,
  не множаться в квадраті тощо). Точніше, лінійне рівняння другого
  порядку з двома незалежними змінними є рівнянням вигляду
\end{enumerate}

\begin{verbatim}
(1.1) Au_{xx} + Bu_{xy} + Cu_{yy} + Du_x + Eu_y + Fu = G
\end{verbatim}

де A, B, C, D, E, F і G --- це константи або задані функції
пояснювальних змінних x і y.

\begin{verbatim}
u_{tt} = e^{-t}u_{xx} + \sin t
\end{verbatim}

(лінійне рівняння), \texttt{uu\_\{xx\}\ +\ u\_t\ =\ 0} (нелінійне
рівняння), \texttt{u\_\{xx\}\ +\ yu\_\{yy\}\ =\ 0} (лінійне рівняння),
\texttt{xu\_x\ +\ yu\_y\ +\ u\^{}2\ =\ 0} (нелінійне рівняння).

\begin{itemize}
\tightlist
\item
  \begin{enumerate}
  \def\labelenumi{\arabic{enumi}.}
  \setcounter{enumi}{3}
  \tightlist
  \item
    Однорідність. Рівняння (1.1) називають однорідним, якщо права
    частина \texttt{G(x,\ y)} ідентична нулю для всіх x і y. Якщо
    \texttt{G(x,\ y)} не ідентичне нулю, то рівняння називають
    неоднорідним.
  \end{enumerate}
\item
  \begin{enumerate}
  \def\labelenumi{\arabic{enumi}.}
  \setcounter{enumi}{4}
  \tightlist
  \item
    Типи шансів. Якщо коефіцієнти A, B, C, D, E і F рівняння (1.1) є
    сталими, то це рівняння називається рівнянням зі сталими
    коефіцієнтами (інакше рівнянням зі змінними коефіцієнтами).
  \end{enumerate}
\item
  \begin{enumerate}
  \def\labelenumi{\arabic{enumi}.}
  \setcounter{enumi}{5}
  \tightlist
  \item
    Три основні типи лінійних рівнянь. Усі лінійні рівняння з частковими
    похідними другого порядку виду (1.1) належать до одного з трьох
    типів: а) параболічних, б) гіперболічних, в) еліптичних.
  \end{enumerate}
\end{itemize}

Параболічний тип. Рівняння параболічного типу описують процеси
теплопровідності та дифузії і визначаються умовою
\texttt{B\^{}2\ -\ 4AC\ =\ 0} .

Гіперболічний тип. Рівняння гіперболічного типу описів

осцилюючі системи та хвильові рухи визначаються умовою
\texttt{B\^{}2-4AC\textgreater{}0} .

Еліпічний тип. Рівняння еліптичного типу описують усталені процеси і
визначаються умовним \texttt{B\^{}2-4AC\textless{}0} . Приклади

\begin{enumerate}
\def\labelenumi{\alph{enumi})}
\tightlist
\item
\end{enumerate}

\begin{verbatim}
u_i = u_{xx}
\end{verbatim}

, \texttt{B\^{}2\ -\ 4AC\ =\ 0} (параболічна), (гіперболічна),
(гіперболічна), (гіперболічна \texttt{y\ \textgreater{}\ 0} ),
(гіперболічна), (еліптична), (еліптична), параболічна на
\texttt{y\ =\ 0} , гіперболічна на \texttt{y\ \textless{}\ 0} .

(У випадку змінних коефіцієнтів тип рівняння може змінюватися від точки
до точки.)

\section{ЗАУВАЖЕННЯ}\label{ux437ux430ux443ux432ux430ux436ux435ux43dux43dux44f}

\begin{itemize}
\tightlist
\item
  \begin{enumerate}
  \def\labelenumi{\arabic{enumi}.}
  \tightlist
  \item
    Загалом, значення \texttt{B\^{}2-4AC} є функцією незалежних змінних.
    Відповідно, базовий тип рівняння може змінюватися в межах визначення
    рівняння, хоча це не є обов'язковим.
  \end{enumerate}
\item
  \begin{enumerate}
  \def\labelenumi{\arabic{enumi}.}
  \setcounter{enumi}{1}
  \tightlist
  \item
    У рівнянні (1.1) незалежними змінними є x і y. У багатьох задачах
    однією з двох незалежних змінних є час, і рівняння (1.1) можна
    записати через x і t.
  \end{enumerate}
\end{itemize}

\pandocbounded{\includegraphics[keepaspectratio]{images/_page_10_Figure_9.jpeg}}

PHC 1.1. Діаграма класифікації диференціальних рівнянь у похідних
похідних.

\begin{enumerate}
\def\labelenumi{\arabic{enumi}.}
\setcounter{enumi}{2}
\tightlist
\item
  Рис. 1.1. Представлено діаграму класифікації для диференціальних
  рівнянь у похідних похідних.
\end{enumerate}

\section{ЗАВДАННЯ}\label{ux437ux430ux432ux434ux430ux43dux43dux44f}

\begin{itemize}
\tightlist
\item
  \begin{enumerate}
  \def\labelenumi{\arabic{enumi}.}
  \tightlist
  \item
    Класифікуємо наступні рівняння відповідно до всіх характеристик,
    зазначених у днаграмі на рис. 1.1:
  \end{enumerate}
\item
  \begin{enumerate}
  \def\labelenumi{\alph{enumi})}
  \tightlist
  \item
    \texttt{u\_t\ =\ u\_\{xx\}\ +\ 2u\_x\ +\ u} ,
  \end{enumerate}
\end{itemize}

\texttt{6)\ \textquotesingle{}u\_t\ =\ u\_\{xx\}\ +\ e\^{}\{-t\}\textquotesingle{}\ ,\ B)\ \textquotesingle{}u\_\{xx\}\ +\ 3u\_\{xy\}\ +\ u\_\{yy\}\ =\ \textbackslash{}sin\ x\textquotesingle{}\ ,\ \textquotesingle{}u\_\{tt\}\ =\ uu\_\{xxx\}\ +\ e\^{}\{-t\}\textquotesingle{}\ .}

\texttt{-\ 2.\ Скільки\ розв\textquotesingle{}язків\ \textquotesingle{}u\_1\ =\ u\_\{xx\}\textquotesingle{}\ рівняння\ існує?\ Спробуйте\ знайти\ рішення,\ як\ \textquotesingle{}u(x,\ t)\ =\ e\^{}\{a\textbackslash{}lambda\ +\ bt\}\textquotesingle{}\ .}
- 3. Якщо функції \texttt{u\_1(x,\ t)} і \texttt{u\_2(x,\ t)}
задовольняють рівняння (1.1), чи задовольняє сума цих функцій його?
Доведи це. - 4. Ймовірно, найпростіше диференціальне рівняння в похідних
---

\begin{verbatim}
\frac{\partial u(x, y)}{\partial x} = 0.
\end{verbatim}

Чи можете ви розв'язати це рівняння? (Знайдіть усі функції, що
задовольняють це рівняння.)

\begin{enumerate}
\def\labelenumi{\arabic{enumi}.}
\setcounter{enumi}{4}
\tightlist
\item
  Що можна сказати про рівняння
\end{enumerate}

\begin{verbatim}
\frac{\partial^2 u(x, y)}{\partial x \partial y} = 0
\end{verbatim}

?

Чи можете ви знайти всі розв'язки цього рівняння? Скільки їх? Порівняйте
з кількістю розв'язків звичайного диференціального рівняння

\begin{verbatim}
\frac{d^2y}{dx^2} = 0.
\end{verbatim}

\section{ЗАДАЧІ ТИПУ ДИФУЗІЇ (ПАРАБОЛІЧНІ
РІВНЯННЯ)}\label{ux437ux430ux434ux430ux447ux456-ux442ux438ux43fux443-ux434ux438ux444ux443ux437ux456ux457-ux43fux430ux440ux430ux431ux43eux43bux456ux447ux43dux456-ux440ux456ux432ux43dux44fux43dux43dux44f}

МЕТА ЛЕКЦІЇ: Показати, як параболічні рівняння використовуються для
розв'язання задач теплопровідності та дифузії. Наведено кілька прикладів
параболічних рівнянь, а також пояснюється фізичне значення різних членів
\texttt{(u\_1,\ u\_2,\ u\_3,\ u\_4,\ u\_6)} , включених у ці рівняння.

На конкретному прикладі розглядається основна ідея зміщеної задачі для
параболічного рівняння. Одна з головних цілей цієї лекції --- надати
читачеві інтуїтивне уявлення

Розуміння параболічних проблем.

Почнемо з розгляньмо просту фізичну задачу і те, як її можна описати
математичною моделлю з рівнянням у похідних похідних.

Потім ми ускладнимо задачу і покажемо, як нові фізичні умови призводять
до появи нових рівнянь у похідних похідних. У цій лекції диференціальні
рівняння не виводяться і не розв'язуються. Про це ми поговоримо пізніше.

\section{Простий експеримент з
теплопровідністю}\label{ux43fux440ux43eux441ux442ux438ux439-ux435ux43aux441ux43fux435ux440ux438ux43cux435ux43dux442-ux437-ux442ux435ux43fux43bux43eux43fux440ux43eux432ux456ux434ux43dux456ux441ux442ux44e}

Припустимо, що ми провели простий експеримент, який розбили на такі
кроки:

КРОК 1. Візьміть достатньо довгий (скажімо,
\texttt{L=2\textbackslash{},\textbackslash{}mathrm\{m\}} ) стрижень
(наприклад, мідний) діаметром
\texttt{2\textbackslash{},\textbackslash{}mathrm\{cm\}} , бічна поверхня
(але не кінці) якої термічно ізольована. Замість стрижня можна
використовувати мідну трубку, покриту теплоізоляцією зовні та зсередини.

Іншими словами, тепло може проходити лише через кінці стрижня і не може
проходити через його бічну поверхню.

КРОК 2. Тепер помістимо цей стрижень у пристрій із фіксованою
температурою \texttt{T\_\{\textbackslash{}mathfrak\{o\}\}} (градуси
Цельсія) на достатньо довгий час, щоб температура всередині стрижня
стала

Так само, як і температура пристрою. Для простоти припустимо, що
температура пристрою \texttt{T\_0\ =\ 10\^{}\{\textbackslash{}circ\}} C.

КРОК 3. Давайте витягнемо стрижень з нашого пристрою у певний момент
часу, що зручно взяти як нуль t = 0, і приєднаємо до нього дві термопари
з кінців. Завдання цих елементів --- підтримувати фіксовані температури
\texttt{T\_1} і \texttt{T\_2} на кінцях (наприклад,
\texttt{T\_1=0\^{}\{\textbackslash{}circ\}} C і
\texttt{T\_2=50\^{}\{\textbackslash{}circ\}} C). Іншими словами, два
термостати постійно контролюють температуру на кінцях стрижня, і якщо
температура відхиляється від встановлених значень \texttt{T\_1} і
\texttt{T\_2} , активуються потужні нагрівальні або охолоджувальні
елементи (пристрої) для відповідного регулювання температури. Наш
експеримент ілюструє RPS, 2 1,

\pandocbounded{\includegraphics[keepaspectratio]{images/_page_13_Figure_4.jpeg}}

Рис. 2.1. Схема експерименту: a --- це термопара, яка підтримує
встановлену температуру \texttt{T\_1} на лівому кінці стрижня;
\texttt{\textbackslash{}delta} --- це термопара, яка підтримує
встановлену температуру \texttt{T\_2} на правому кінці стрижня.

КРОК 4. Тепер ми будемо слідкувати за температурним профілем у стрижні
за допомогою пристрою візуалізації. Чому ми хочемо провести такий
експеримент? Це інше питання, про яке ми поговоримо пізніше.

Це кінець опису нашого експерименту.

Головна мета цієї лекції --- показати, як цю фізичну проблему та її
модифікації можна змоделювати мовою параболічних рівнянь.

\section{Математична модель
теплопровідності}\label{ux43cux430ux442ux435ux43cux430ux442ux438ux447ux43dux430-ux43cux43eux434ux435ux43bux44c-ux442ux435ux43fux43bux43eux43fux440ux43eux432ux456ux434ux43dux43eux441ux442ux456}

Опис нашої фізичної проблеми вимагає трьох типів відносин.

\begin{itemize}
\tightlist
\item
  \begin{enumerate}
  \def\labelenumi{\arabic{enumi}.}
  \tightlist
  \item
    Рівняння з частковими похідними, які описують фізичне явище
    теплопровідності.
  \end{enumerate}
\item
  \begin{enumerate}
  \def\labelenumi{\arabic{enumi}.}
  \setcounter{enumi}{1}
  \tightlist
  \item
    Граничні умови, що описують процеси теплопередачі на межах.
  \end{enumerate}
\item
  \begin{enumerate}
  \def\labelenumi{\arabic{enumi}.}
  \setcounter{enumi}{2}
  \tightlist
  \item
    Початкові твердження, які описують стан системи на початку процесу.
  \end{enumerate}
\end{itemize}

\section{Рівняння
теплопровідності}\label{ux440ux456ux432ux43dux44fux43dux43dux44f-ux442ux435ux43fux43bux43eux43fux440ux43eux432ux456ux434ux43dux43eux441ux442ux456}

Базове одновимірне рівняння теплопровідності записується у вилці

(2.1)

\begin{verbatim}
u_t = \alpha^2 u_{xx}, \quad 0 < x < L, \quad 0 < t < \infty.
\end{verbatim}

Це рівняння пов'язує значення \texttt{u\_t} --- швидкість зміни
температури з часом (виміряна в градусах/с)

та величину \texttt{u\_\{xx\}} --- це увігнутий профілю температури
\texttt{u(x,\ t)} (який є мірою різниці між температурою в заданій точці
та температурою в сусідніх точках).

\pandocbounded{\includegraphics[keepaspectratio]{images/_page_14_Figure_14.jpeg}}

Рис. 2.2. Стрілки показують зміну температури відповідно до рівняння
\texttt{u\_t\ =\ \textbackslash{}alpha\^{}2\ u\_\{xx\}} .

\texttt{Це\ рівняння\ буде\ виведено\ з\ базового\ рівняння\ збереження\ тепла\ в\ наступних\ лекціях,\ але\ наразі\ ми\ розглянемо\ його\ заданим.\ Це\ рівняння\ каже,\ що\ температура\ \textquotesingle{}u\textbackslash{}left(x,\textbackslash{}\ t\textbackslash{}right)\textquotesingle{}\ (у\ певній\ точці\ t\ у\ певній\ точці\ x\ стрижня)\ зростає\ \textquotesingle{}(u\_t\textgreater{}0)\textquotesingle{}\ або\ зменшується\ \textquotesingle{}(u\_t\textless{}0)\textquotesingle{}\ залежно\ від\ того,\ чи\ є\ друга\ похідна\ \textquotesingle{}u\_\{xx\}\textquotesingle{}\ позитивною\ чи\ негативною.\ Фіг.\ Рисунок\ 2.2\ ілюструє\ зміну\ температури\ в\ різних\ точках\ на\ елементі.}

Давайте подивимось, як значення \texttt{u\_\{xx\}} можна інтерпретувати
мовою теплопровідності. Припустимо, ми наблизимо значення
\texttt{u\_\{xx\}} фінальних постів

\begin{verbatim}
u_{xx}(x, t) \cong \frac{1}{\Delta x^2} [u(x + \Delta x, t) - 2u(x, t) + u(x - \Delta x, t)].
\end{verbatim}

Це співвідношення можна переписати у вигляді

\begin{verbatim}
u_{\lambda x}(x, t) \simeq \frac{2}{\Delta x^2} \left[ \frac{u(x + \Delta x, t) + u(x - \Delta x, t)}{2} - u(x, t) \right].
\end{verbatim}

Тепер ми можемо дати наступну інтерпретацію значення \texttt{u\_\{xx\}}
.

\begin{enumerate}
\def\labelenumi{\arabic{enumi}.}
\item
  Якщо температура \texttt{u(x,\ t)} менша за середнє значення
  температури у двох сусідніх точках, то
  \texttt{u\_\{xx\}\ \textgreater{}\ 0} (тут загальний тепловий потік по
  осі x додатний).
\item
  Якщо температура \texttt{u(x,\ t)} дорівнює середньому значенню
  температур у двох сусідніх точках, то \texttt{u\_\{xx\}\ =\ 0} (тут
  загальний тепловий потік по осі x дорівнює нулю).
\end{enumerate}

\texttt{3.\ Якщо\ температура\ \textquotesingle{}u\textbackslash{}left(x,\textbackslash{}\ t\textbackslash{}right)\textquotesingle{}\ вища\ за\ середню\ температуру\ у\ двох\ сусідніх\ точках,\ то\ \textquotesingle{}u\_\{xx\}\textless{}0\textquotesingle{}\ (тут\ —\ повний\ потік}

тепла вздовж осі x є від'ємною).

\texttt{Усе\ це\ ілюструється\ на\ рис.\ 2.2.\ Іншими\ словами,\ якщо\ температура\ в\ точці\ x\ більша\ за\ середню\ температуру\ у\ двох\ сусідніх\ точках\ \textquotesingle{}x\ -\ \textbackslash{}Delta\ x\textquotesingle{}\ і\ \textquotesingle{}x\ +\ \textbackslash{}Delta\ x\textquotesingle{}\ ,\ то\ температура\ в\ точці\ x\ зменшиться.\ Відповідно,\ точна\ швидкість\ зниження\ температури,\ а\ саме\ величина\ \textquotesingle{}u\_t\textquotesingle{}\ ,\ пропорційна\ цій\ різниці.\ Коефіцієнт\ пропорційності\ \textquotesingle{}\textbackslash{}alpha\^{}2\textquotesingle{}\ визначається\ властивостями\ матеріалу.\ Ми\ розглянемо\ цю\ константу\ детальніше\ у\ майбутніх\ лекціях.}

\section{Граничні умови
(ГІ)}\label{ux433ux440ux430ux43dux438ux447ux43dux456-ux443ux43cux43eux432ux438-ux433ux456}

Усі фізичні проблеми характеризуються наявністю певних меж. Отже, щоб
рухатися вперед, ми повинні включити межі в математичну модель, щоб
адекватно описати фізичну проблему. У нашому експерименті граничні умови
(ГІ) дуже прості. Оскільки температура на кінцях x=0 і x=L завжди
залишається сталою і дорівнює \texttt{T\_1} та \texttt{T\_2} , можливо
записати

(2.2)

\begin{verbatim}
\left\{ \begin{array}{l} u\left(0,\ t\right) = T_{t}, \\ u\left(L,\ t\right) = T_{2}, \end{array} \right. \quad 0 < t < \infty,
\end{verbatim}

\section{Початкові умови
(NU)}\label{ux43fux43eux447ux430ux442ux43aux43eux432ux456-ux443ux43cux43eux432ux438-nu}

Усі фізичні процеси мають починатися у певний момент часу (зазвичай
приймається як нуль t=0), тому ми повинні встановити фізичні умови саме
в цю точку. Відтоді, як ми почали контролювати температуру з моменту,
коли паличка мала постійну температуру
\texttt{T\_\{\textbackslash{}mathrm\{\ \}o\}} , ми

(2.3)

\begin{verbatim}
u(x, 0) = T_0, \quad 0 \le x \le L
\end{verbatim}

Тепер ми створили математичну модель експерименту. Система співвідношень
(2.1), (2.2) та (2.3) називається задачею зміщення для рівняння
теплопровідності і зазвичай записується як

(2.4)

\begin{verbatim}
\begin{aligned} & (\text{YUII}) \quad u_t = \alpha^2 u_{xx}, & 0 < x < L, \ 0 < t < \infty, \\ & \left( \Gamma \text{Y} \right) \quad \begin{cases} u \ (0, \ t) = T_1, \\ u \ (L, \ t) = T_2, \end{cases} & 0 < t < \infty, \\ & (\text{HY}) \quad u \ (x, \ 0) = T_0, \quad 0 \leqslant x \leqslant L. \end{aligned}
\end{verbatim}

Цікаво, що це зовсім не очевидно: чи існує одна функція
\texttt{u(x,\ t)}, яка задовольняє умови (2.4), і чи описує ця функція
температуру в стрижні? Отже, наше безпосереднє завдання --- знайти цю
функцію \texttt{u(x,\ t)} з (2.4).

Перш ніж завершити цю лекцію, давайте розглянемо деякі варіації основної
проблеми. Почнемо з невеликих змін рівняння теплопровідності.

\section{Деякі рівняння дифузійного
типу}\label{ux434ux435ux44fux43aux456-ux440ux456ux432ux43dux44fux43dux43dux44f-ux434ux438ux444ux443ux437ux456ux439ux43dux43eux433ux43e-ux442ux438ux43fux443}

Теплопередача через бічну поверхню пропорційна різниці температур

Рівняння

\begin{verbatim}
u_t = \alpha^2 u_{xx} - \beta (u - u_0)
\end{verbatim}

Описує теплопровідність елемента, враховуючи не лише дифузію теплових
\texttt{\textbackslash{}alpha\^{}2\ u\_\{xx\}} вздовж елемента, а й
теплопередачу через бічну поверхню елемента. Відтік тепла (
\texttt{\textbackslash{}beta\ \textgreater{}\ 0} ) або його приплив (
\texttt{\textbackslash{}beta\ \textless{}\ 0} ) пропорційний різниці між
температурою u стрижня та температурою навколишнього середовища
\texttt{u\_0} ( \texttt{\textbackslash{}beta} --- це сталий коефіцієнт
пропорційності). Якщо коефіцієнт \texttt{\textbackslash{}beta} великий
порівняно з \texttt{\textbackslash{}alpha\^{}2} , то теплопотік уздовж
стрижня буде малим порівняно з

Протікання через бічну поверхню, і тому тепло буде проходити через неї
відповідно до закону
\texttt{u\_t\ =\ -\textbackslash{}beta\ (u\ -\ u\_a)} .

У хімії, де величина відіграє роль концентрації, рівняння

\begin{verbatim}
u_1 = \alpha^2 u_{xx} - \beta (u - u_0)
\end{verbatim}

означає, що швидкість зміни \texttt{u\_i} кількості речовини залежить як
від дифузії \texttt{\textbackslash{}alpha\^{}2\ u\_\{xx\}} (у напрямку
осі x), так і від появи ( \texttt{\textbackslash{}beta\ \textless{}\ 0}
) або розпаду ( \texttt{\textbackslash{}beta\ \textgreater{}\ 0} )
речовини в хімічній реакції, і пропорційна різниці між двома
концентраціями u і \texttt{u\_\{\textbackslash{}theta\}} .

\section{Внутрішнє джерело
тепла}\label{ux432ux43dux443ux442ux440ux456ux448ux43dux454-ux434ux436ux435ux440ux435ux43bux43e-ux442ux435ux43fux43bux430}

Неоднородне рівняння

\begin{verbatim}
u_t = \alpha^2 u_{xx} + f(x, t)
\end{verbatim}

відповідає випадку, коли всередині елемента є джерело тепла (розташоване
вздовж усього елемента і діє в будь-який момент t). Наприклад, може бути
так, що дріт з електричним струмом проходить через стрижень і опір є
постійним джерелом тепла \texttt{f(x,t)} = K.

\section{Рівняння конвективної
дифузії}\label{ux440ux456ux432ux43dux44fux43dux43dux44f-ux43aux43eux43dux432ux435ux43aux442ux438ux432ux43dux43eux457-ux434ux438ux444ux443ux437ux456ux457}

\texttt{Припустимо,\ що\ домішка\ поширюється\ вздовж\ потоку,\ що\ рухається\ зі\ швидкістю\ v.\ Очевидно,\ що\ концентрація\ \textquotesingle{}u\textbackslash{}left(x,\textbackslash{},t\textbackslash{}right)\textquotesingle{}\ домішок\ змінюється\ залежно\ від\ x\ (позитивний\ напрямок\ осі\ x\ обирається\ вздовж\ потоку)\ і\ часу\ t.\ Швидкість\ зміни\ концентрації\ \textquotesingle{}u\_t\textquotesingle{}\ визначається\ рівнянням\ конвективної\ дифузії}

\begin{verbatim}
u_t = \alpha^2 u_{xx} - v u_x.
\end{verbatim}

Термін \texttt{\textbackslash{}alpha\^{}2\ u\_\{xx\}} описує внесок
дифузії, а \texttt{vu\_x} --- конвективний компонент. Що важливіше,
дифузія або конвекція залежить від співвідношення
\texttt{\textbackslash{}alpha\^{}2} до v. Димові частинки піднімаються
конвективно через рух гарячого повітря і одночасно розсіюють через
вихорові рухи повітря. Окрім цих модифікацій рівняння теплопровідності,
можливо змінювати граничні умови на кінцях стрижня відповідно до
фізичної ситуації. Деякі з цих змін ми обговоримо у третій лекції.

\section{ЗАУВАЖЕННЯ}\label{ux437ux430ux443ux432ux430ux436ux435ux43dux43dux44f-1}

Рівняння теплопровідності
\texttt{u\_t\ =\ \textbackslash{}alpha\^{}2(x)\ u\_\{xx\}} з змінним
коефіцієнтом \texttt{\textbackslash{}alpha(x)} відповідає задачі, в якій
дифузія залежить від x (тобто матеріал є гетерогенним). Наприклад, мідь
і сталь---

\pandocbounded{\includegraphics[keepaspectratio]{images/_page_18_Picture_4.jpeg}}

PHC. 2.3.

Пластини контактують, як показано на рис. 2.3. Припустимо, що
температура \texttt{u(0,t)=0\^{}\{\textbackslash{}circ\}} C
підтримується на лівому боці міді, а температура
\texttt{u(L,t)=20\^{}\{\textbackslash{}circ\}} C --- з правого боку
сталі, тоді рівняння, що описує теплопровідність у цій системі, виглядає
так:

\begin{verbatim}
u_t = \alpha^2(x) \, u_{_{\! A \, X}}, \quad 0 < x < L,
\end{verbatim}

де

\begin{verbatim}
\alpha(x) = \begin{cases}
\alpha_1, & \text{коефіцієнт дифузії в міді}, \quad 0 < x < L/2, \\
\alpha_2, & \text{коефіцієнт дифузії в сталі}, \quad L/2 < x < L.
\end{cases}
\end{verbatim}

\section{ЗАВДАННЯ}\label{ux437ux430ux432ux434ux430ux43dux43dux44f-1}

\begin{enumerate}
\def\labelenumi{\arabic{enumi}.}
\tightlist
\item
  Якщо початкова температура в стрижні
\end{enumerate}

\begin{verbatim}
u(x, 0) = \sin \pi x, \quad 0 \leqslant x \leqslant 1,
\end{verbatim}

а граничні умови такі:

\begin{verbatim}
u(0, t) = 0, \qquad u(1, t) = 0.
\end{verbatim}

Тоді як поводитиметься температура в стрижні \texttt{U(x,\ t)} при
\texttt{t\ \textgreater{}\ 0} ? 2. Припустимо, що в стрижні є постійне
внутрішнє джерело тепла, так що рівняння теплопровідності в стрижні
записується як

\begin{verbatim}
u_t = \alpha^2 u_{xx} + 1, \quad 0 < x < 1.
\end{verbatim}

Припустимо, що ми зафіксували температуру на межі з значеннями
\texttt{u(0,\ t)} = 0 і \texttt{u(1,\ t)} = 1. Якою буде стаціонарна
температура у стрижні? Іншими словами, який c-незалежний розподіл
\texttt{U(x,\ t)} перешкоджує \texttt{u(x,\ t)}?

\begin{enumerate}
\def\labelenumi{\arabic{enumi}.}
\setcounter{enumi}{2}
\tightlist
\item
  Припустимо, що металевий стрижень втрачає тепло через бічну поверхню
  згідно з рівнянням
\end{enumerate}

\begin{verbatim}
u_1 = \alpha^2 u_{xx} - \beta u
\end{verbatim}

, \texttt{0\ \textless{}\ x\ \textless{}\ 1} .

І припустимо, що ми підтримуємо кінці стрижня при температурах
\texttt{U(0,\ T)} = 1, \texttt{U(1,\ T)} = 1. Знайдіть стаціонарний
розподіл температури по стрижню. Побудуйте графік цього розподілу. У
яких моментах відбувається тепловий потік?

\texttt{4.\ Припустимо,\ що\ термоізольований\ металевий\ стрижень\ довжини\ L\ =\ 1\ має\ початкову\ температуру\ \textquotesingle{}\textbackslash{}sin(3\textbackslash{}pi\ x)\textquotesingle{}\ ,\ а\ фіксовану\ температуру\ —\ 0\ і\ \textquotesingle{}10\textbackslash{},\^{}\{\textbackslash{}circ\}\textquotesingle{}\ C\ на\ лівому\ та\ правому\ кінцях.}

\section{Лекція 3}\label{ux43bux435ux43aux446ux456ux44f-3}

\section{ГРАНИЧНІ УМОВИ В ЗАДАЧАХ ДИФУЗІЙНОГО
ТИПУ}\label{ux433ux440ux430ux43dux438ux447ux43dux456-ux443ux43cux43eux432ux438-ux432-ux437ux430ux434ux430ux447ux430ux445-ux434ux438ux444ux443ux437ux456ux439ux43dux43eux433ux43e-ux442ux438ux43fux443}

МЕТА ЛЕКЦІЇ: Показати, як виникають різні граничні умови у задачах
теплопровідності та дифузії, а також ввести важливе поняття теплового
потоку.

Розглядаються три найважливіші типи граничних умов:

\begin{itemize}
\tightlist
\item
  \begin{enumerate}
  \def\labelenumi{\arabic{enumi}.}
  \tightlist
  \item
    u = \texttt{g(t)} (температура задається на межі),
    \texttt{-\ \textquotesingle{}2.\ \textbackslash{}\ \textbackslash{}frac\{\textbackslash{}partial\ u\}\{\textbackslash{}partial\ n\}\ +\ \textbackslash{}lambda\ u\ =\ g\textbackslash{}left(t\textbackslash{}right)\textquotesingle{}\ (встановлення\ температури\ навколишнього\ середовища);}
  \end{enumerate}
\end{itemize}

n --- вектор зовнішньої нормалі до межі),
\texttt{3.\ \textbackslash{}frac\{\textbackslash{}partial\ u\}\{\textbackslash{}partial\ n\}\ =\ g(t)}
(визначений тепловий потік через межу).

У задачах теплопровідності зазвичай використовують три основні типи
граничних умов. Ця лекція наводить приклади того, як різні фізичні умови
експерименту призводять до різних типів граничних умов.

\section{Граничні умови першого виду (температура задається на
межі)}\label{ux433ux440ux430ux43dux438ux447ux43dux456-ux443ux43cux43eux432ux438-ux43fux435ux440ux448ux43eux433ux43e-ux432ux438ux434ux443-ux442ux435ux43cux43fux435ux440ux430ux442ux443ux440ux430-ux437ux430ux434ux430ux454ux442ux44cux441ux44f-ux43dux430-ux43cux435ux436ux456}

Розглянемо тепловий потік в одномірному стрижні, показаний на рис. 3 1,
і припустимо, що температура кінців змінюється відповідно до законів
\texttt{g\_1(t)} і \texttt{g\_2(t)} .

\pandocbounded{\includegraphics[keepaspectratio]{images/_page_20_Picture_4.jpeg}}

Фіг. 31. Температура буде відновлена на межі.

Як згадувалося в попередній лекції, кожен пристрій, що підтримує
встановлену температуру на кінцях, складається з термостата та
нагрівального елемента для підтримки відповідних теплових умов. Проблеми
з граничними умовами першого виду дуже поширені. У деяких випадках
складність полягає у пошуку контролерів граничних умов, тобто температур
\texttt{g\_1(t)} та \texttt{g\_2(t)} меж, які призводять до певної зміни
температури всередині стрижня. Наприклад, у металургії часто необхідно
обирати граничні умови контролю, щоб температура металу всередині печі
змінювалася з часом, а температурний градієнт був невеликим.

\pandocbounded{\includegraphics[keepaspectratio]{images/_page_20_Picture_7.jpeg}}

Фіг. 3.2. Зниження температури на межі.

Граничні умови першого виду також виникають у багатовимірних задачах
теплопровідності. Як приклад можна навести цікаву задачу визначення
температурного поля всередині кругового диска з радіусом \emph{R} (див.
рис. 3.2), коли гранична температура задається у полярних координатах
формулою
\texttt{u(R,\ \textbackslash{}theta,\ t)\ =\ \textbackslash{}cos\ t\ \textbackslash{}sin\ \textbackslash{}theta}
.

Звісно, щоб розв'язати задачу, потрібно знати початкову температуру, але
в цьому випадку початкові умови перестають впливати на розв'язок через
дуже короткий час, і отримана температура всередині диска визначається
лише значеннями температури на межі.

\section{Граничні умови другого виду (задана температура навколишнього
середовища)}\label{ux433ux440ux430ux43dux438ux447ux43dux456-ux443ux43cux43eux432ux438-ux434ux440ux443ux433ux43eux433ux43e-ux432ux438ux434ux443-ux437ux430ux434ux430ux43dux430-ux442ux435ux43cux43fux435ux440ux430ux442ux443ux440ux430-ux43dux430ux432ux43aux43eux43bux438ux448ux43dux44cux43eux433ux43e-ux441ux435ux440ux435ux434ux43eux432ux438ux449ux430}

Припустимо, що ми знову розглядаємо термічно ізольований мідний
стержень, але тепер відмовляємося від заданого температурного режиму на
кінцях стрижня. Натомість давайте наблизимо кінці стрижня до контакту з
двома середовищами. Нехай температура одного з них \texttt{g\_1(t)} , а
іншого \texttt{g\_2(t)} . Іншими словами, припустимо, що лівий кінець
стрижня поміщений у контейнер із рідиною, температура якої змінюється
відповідно до закону \texttt{g\_1(t)} . А правий кінець поміщається в
іншу рідину з температурою \texttt{g\_2(t)} (рис. 3, 3)

\pandocbounded{\includegraphics[keepaspectratio]{images/_page_21_Picture_6.jpeg}}

Рис. 3.3. Конвективна теплопередача через межі, a --- це рідина при
температурі \texttt{g\_1(t)} , \texttt{\textbackslash{}delta} --- рідина
при температурі \texttt{g\_2(t)} .

Задаючи граничні умови такого типу, ми не можемо вважати, що граничні
температури стрижня збігаються з температурами рідини. Але ми знаємо
закон Ньютона: якщо температура одного з кінців стрижня менша за
температуру відповідної рідини, тепло потече в стрижень зі швидкістю,
пропорційною різницям температур. Іншими словами, для одномірного
стрижня з межами x=0 і x=1 закон теплопередачі Ньютона формулюється так:

(3.1)

\begin{verbatim}
\{
\end{verbatim}

\texttt{вихідний\ тепловий\ потік\ на\ \textquotesingle{}x=0\textquotesingle{}\ дорівнює\ \textquotesingle{}h\textbackslash{}left{[}u\textbackslash{}left(0,\textbackslash{},t\textbackslash{}right)-g\_\{1\}\textbackslash{}left(t\textbackslash{}right)\textbackslash{}right{]}\textquotesingle{}\ ,\ вихідний\ тепловий\ потік\ при\ \textquotesingle{}x=L\textquotesingle{}\ дорівнює\ \textquotesingle{}h\textbackslash{}left{[}u\textbackslash{}left(L,\textbackslash{},t\textbackslash{}right)-g\_\{2\}\textbackslash{}left(t\textbackslash{}right)\textbackslash{}right{]}\textquotesingle{}\ ,}

де h --- коефіцієнт теплопередачі, який показує, скільки калорій
проходить через межу за одну секунду при різниці температур у один
градус. Вихідний потік дорівнює кількості кало-

Пройшов крізь палочний спис за одну секунду. Зверніть увагу, що вихідний
тепловий потік буде позитивним на кінці стрижня, де температура стрижня
вища за температуру навколишнього середовища. Рівняння (3.1) разом із
відомим законом теплопровідності Фур'є тепер можна використовувати для
отримання Р и с

\pandocbounded{\includegraphics[keepaspectratio]{images/_page_22_Figure_3.jpeg}}

. 3.4. Ілюстрація закону Фур'є:
\texttt{\textbackslash{}frac\{\textbackslash{}partial\ u\}\{\textbackslash{}partial\ n\}\ \textgreater{}\ 0}
означає, що тепло надходить у цю область;
\texttt{\textbackslash{}frac\{\textbackslash{}partial\ u\}\{\textbackslash{}partial\ n\}\ \textless{}\ 0}
означає, що тепло виходить із цієї зони; n --- напрямок зовнішньої
нормалі;
\texttt{\textbackslash{}frac\{\textbackslash{}partial\ u\}\{\textbackslash{}partial\ n\}\ -}
Зміна температури в напрямку N;
\texttt{\textbackslash{}frac\{\textbackslash{}partial\ u\}\{\textbackslash{}partial\ n\}\ =\ -\textbackslash{}frac\{\textbackslash{}partial\ u\}\{\textbackslash{}partial\ (-n)\}}
є похідною у напрямку внутрішньої нормальної --- це похідна, взята
\texttt{\textbackslash{}mathbf\{e\}} протилежному знаку протилежного
знака пормалі.

граничні умови. Закон Фур'є дає друге представлення (перше міститься у
(3.1)) для вихідного теплового потоку. Прирівнюючи ці два вирази до
потоку, ми отримуємо граничні умови, які шукаємо. Сформулюємо закон
Фур'є (встановлений експериментально):

(3.2) Тепловий потік, що проходить через межу площі, пропорційний
нормальній похідній температури у напрямку внутрішньої норми.

Цей закон стверджує, що якщо температура швидко підвищується у напрямку
зовнішньої температури до межі D (рисунок 3.4), то потік буде текти з
навколишнього середовища до області D.

У нашій одновимірній задачі закон Фур'є має вигляд:

(3.3)

\begin{verbatim}
\begin{cases} \text{вытекающий поток тепла при } x=0 \text{ равен } k \frac{\partial u}{\partial x}, \\ \text{вытекающий поток тепла при } x=' \text{ равен } -k \frac{\partial u}{\partial x}, \end{cases}
\end{verbatim}

де k --- теплопровідність матеріалу, яка є мірою того, наскільки добре
матеріал проводить тепло. Погано провідні матеріали мають коефіцієнт
близький до нуля в одиницях GHS, тоді як мідь та алюміній мають значення
значно вище нуля.

\pandocbounded{\includegraphics[keepaspectratio]{images/_page_23_Figure_5.jpeg}}

Фіг. 3.5. Ще одна ілюстрація закону Чотирьох.

Закон Фур'є (3.3) фактично описує теплопровідність всюди всередині
стрижня, а не лише на його межі; Наприклад, (див. рис. 3.5).

(3.4) Потік, що проходить через точку

x

\texttt{зліва\ направо\ \textquotesingle{}=-k\textbackslash{}frac\{\textbackslash{}partial\ n\}\{\textbackslash{}partial\ x\}\textquotesingle{}\ Закон\ Фур\textquotesingle{}є\ (3.4)\ стверджує:\ якщо\ \textquotesingle{}u\_x\ \textless{}\ 0\textquotesingle{}\ ,\ то\ в\ точці\ \textquotesingle{}x\_0\textquotesingle{}\ тепло\ тече\ зліва\ направо,\ якщо\ \textquotesingle{}u\_x\ \textgreater{}\ 0\textquotesingle{}\ ,\ то\ в\ точці\ \textquotesingle{}x\_0\textquotesingle{}\ тепло\ тече\ справа\ наліво.\ (Тепло\ завжди\ переходить\ від\ вищих\ температур\ до\ нижчих.)}

Нарешті, якщо використовувати вирази (3.1) і (3.3) для теплового потоку,
то отримуємо бажані граничні умови для задачі, описаної на рис. 3.3, у
чисто математичній формі

\begin{verbatim}
(\Gamma Y) \quad \begin{cases} \frac{\partial u\left(0,\,t\right)}{\partial x} = \frac{h}{k}\left[u\left(0,\,t\right) - g_{1}\left(t\right)\right], \\ \frac{\partial u\left(L,\,t\right)}{\partial x} = -\frac{h}{k}\left[u\left(L,\,t\right) - g_{2}\left(t\right)\right]. \end{cases} \quad 0 < t < \infty,
\end{verbatim}

Дуже часто константа h/k позначається як \texttt{\textbackslash{}lambda}
, а граничні умови теплового потоку на межі записуються у вигляді

(3.5)

\begin{verbatim}
u_{x}(0, t) = \lambda [u(0, t) - g_{1}(t)], u_{x}(L, t) = -\lambda [u(L, t) - g_{2}(t)].
\end{verbatim}

Для випадку вищих розмірностей отримують подібні граничні умови.
Наприклад, якщо межа кругового диска осяяна рухомою рідиною з
температурою \texttt{g(\textbackslash{}theta,t)} , то граничні умови
записуються як

\begin{verbatim}
\frac{\partial u}{\partial r}(R, \theta, t) = -\frac{h}{k} [u(R, \theta, t) - g(\theta, t)].
\end{verbatim}

Тут
\texttt{\textbackslash{}frac\{\textbackslash{}partial\ u\}\{\textbackslash{}partial\ r\}(R,\ \textbackslash{}theta,\ t)}
позначає зовнішню нормальну похідну (у позитивному напрямку r-осі)
значення u, обчислену в точці \texttt{(R,\ \textbackslash{}theta)} межі.
Ми назвемо такі граничні умови лінійними (оскільки вони лінійні в U і
\texttt{u\_r} ), але неоднорідними, оскільки функція
\texttt{g(\textbackslash{}theta,\ t)} знаходиться праворуч.

\section{Граничні умови третього виду (враховуючи потік; зокрема,
враховується випадок термічно ізольованих
меж)}\label{ux433ux440ux430ux43dux438ux447ux43dux456-ux443ux43cux43eux432ux438-ux442ux440ux435ux442ux44cux43eux433ux43e-ux432ux438ux434ux443-ux432ux440ux430ux445ux43eux432ux443ux44eux447ux438-ux43fux43eux442ux456ux43a-ux437ux43eux43aux440ux435ux43cux430-ux432ux440ux430ux445ux43eux432ux443ux454ux442ux44cux441ux44f-ux432ux438ux43fux430ux434ux43eux43a-ux442ux435ux440ux43cux456ux447ux43dux43e-ux456ux437ux43eux43bux44cux43eux432ux430ux43dux438ux445-ux43cux435ux436}

Жоден потік не проходить через термоізольовані межі, тому нормальна
похідна (зовнішня чи внутрішня) має зрівнятися з нулем на межі, оскільки
потік пропорційний нормальній похідній. У випадку одномірного елемента з
ізольованими кінцями x=0 і x=L, граничні умови мають вигляд

\begin{verbatim}
u_{\lambda}(0, t) = 0,
\end{verbatim}

\texttt{\textquotesingle{}u\_\{\textbackslash{}lambda\}(L,\ t)\ =\ 0.\textquotesingle{}\ \textquotesingle{}0\ \textless{}\ t\ \textless{}\ \textbackslash{}infty,\textquotesingle{}\ У\ двовимірній\ області\ ізоляція\ межі\ означає,\ що\ нормальна\ похідна\ температури\ на\ межі\ перетворюється\ на\ нуль.\ Припустимо,\ наприклад,\ що\ круглий\ диск\ термічно\ ізольований\ уздовж\ межі.\ Тоді\ граничну\ умову\ слід\ записати\ так:\ \textquotesingle{}u\_r(R,\ \textbackslash{}theta,\ t)\ =\ 0\textquotesingle{}\ для\ всіх\ \textquotesingle{}0\ \textbackslash{}le\ \textbackslash{}theta\ \textless{}\ 2\textbackslash{}pi\textquotesingle{}\ і\ всіх\ \textquotesingle{}0\ \textless{}\ t\ \textless{}\ \textbackslash{}infty\textquotesingle{}\ .}

З іншого боку, якщо вказати кількість тепла, що проходить через межу
нашого диска, то граничні умови набувають вигляду

\begin{verbatim}
u_r(R, \theta, t) = f(\theta, t),
\end{verbatim}

де \texttt{f(0,\ t)} --- це кількість тепла, що переходить у круглий
диск із зовнішнього джерела.

Тепер проілюструємо різні типи граничних умов.

\section{Типові граничні умови для задачі одномірної
теплопровідності}\label{ux442ux438ux43fux43eux432ux456-ux433ux440ux430ux43dux438ux447ux43dux456-ux443ux43cux43eux432ux438-ux434ux43bux44f-ux437ux430ux434ux430ux447ux456-ux43eux434ux43dux43eux43cux456ux440ux43dux43eux457-ux442ux435ux43fux43bux43eux43fux440ux43eux432ux456ux434ux43dux43eux441ux442ux456}

Припустимо, що у нас є мідний стрижень довжиною 200 см, бікова поверхня
якого термічно ізольована, а початкова температура становить 0 °C.
Припустимо також, що верхній кінець стрижня x=0 термічно ізольований, а
нижній кінець x=200 змивається рухомою водою, яка має постійну
температуру \texttt{g\_2(t)=20} °C (див. рис. 3.6).

\pandocbounded{\includegraphics[keepaspectratio]{images/_page_25_Picture_6.jpeg}}

Фіг. 3.6. Змішане завдання.

Математична модель цієї задачі описується наступними чотирма
відношеннями:

(3.6)

\begin{verbatim}
\begin{aligned}
(\text{УЧП})\quad & u_t = \alpha^2 u_{xx}, && 0 < x < 200,\ 0 < t < \infty, \\
(\text{ГУ})\quad & u_x(0,t) = 0, && 0 < t < \infty, \\
& u_x(200,t) = -\frac{h}{k}\left[u(200,t) - 20\right], && 0 < t < \infty, \\
(\text{НУ})\quad & u(x,0) = 0^\circ\mathrm{C}, && 0 \leqslant x \leqslant 200.
\end{aligned}
\end{verbatim}

де \texttt{\textbackslash{}alpha\^{}2\ =\ 1,16} см 2/с --- теплова
дифузія міді;

k = 0,93 кал/см·с·°C --- теплопровідність міді;

h --- коефіцієнт теплопередачі. Визначення цього співвідношення ---
велика проблема. Він вимірюється швидкістю виходу тепла від другого
кінця стрижня до води. Це залежить від швидкості потоку води, стану
поверхні, яку слід змивати тощо.

Читач може спробувати провести експеримент, щоб визначити

ділення значення цього коефіцієнта.

\section{ЗАУВАЖЕННЯ}\label{ux437ux430ux443ux432ux430ux436ux435ux43dux43dux44f-2}

\begin{enumerate}
\def\labelenumi{\arabic{enumi}.}
\tightlist
\item
  Типова задача внутрішньої теплопровідності для квадрата показана на
  рис. 3.7.
\end{enumerate}

\texttt{У\ цій\ задачі\ початкова\ температура\ \textquotesingle{}u\textbackslash{}left(x,\textbackslash{},y,\textbackslash{},t\textbackslash{}right)\textquotesingle{}\ при\ t\ =\ 0\ задається\ всередині\ квадрата,\ і\ рівняння\ з\ похідними\ з\ частковими\ має\ вигляд\ Р\ и\ с}

\pandocbounded{\includegraphics[keepaspectratio]{images/_page_26_Figure_9.jpeg}}

. 3.7. Типові граничні умови для задачі дифузії в квадраті: a ---
кількість тепла, \texttt{f\_1(t)} проходить через межу; b ---
термоізольована межа; a --- це навколишня температура Равії
\texttt{10\^{}\{\textbackslash{}circ\}} C; z --- температура
підтримується рівною \texttt{10\^{}\{\textbackslash{}circ\}} C.

а граничні умови будуть виконані на
\texttt{0\ \textless{}\ t\ \textless{}\ \textbackslash{}infty} і
визначатимуть температуру в будь-який момент. Якою б не була ця
температура, вона повинна відповідати граничним умовам, зазначеним на
рис. 3.7.

\begin{enumerate}
\def\labelenumi{\arabic{enumi}.}
\setcounter{enumi}{1}
\tightlist
\item
  Зверніть увагу, що гранична умова
\end{enumerate}

\begin{verbatim}
u_r(R, \theta, t) = -\frac{h}{k} [u(R, \theta, t) - g(\theta, t)]
\end{verbatim}

На колі не означає, що гранична температура дорівнює
\texttt{g(\textbackslash{}theta,t)} , але якщо коефіцієнт теплопередачі
h великий, то ці умови фактично еквівалентні вимогі
\texttt{g(\textbackslash{}theta,t)} температури на межі.

\section{ЗАВДАННЯ}\label{ux437ux430ux432ux434ux430ux43dux43dux44f-2}

\begin{enumerate}
\def\labelenumi{\arabic{enumi}.}
\item
  Намалюйте контури графіків для розв'язання змішаної задачі (3.6) для
  різних моментів часу. Чи відповідають ваші графи визначеним граничним
  умовам? Яка стаціонарна температура стрижня? Очевидно, ваша відповідь
  базується на інтуїції?
\item
  Як би ви інтерпретували наступне змішане завдання?
\end{enumerate}

\begin{verbatim}
\begin{aligned}
(\text{УЧП})\quad & u_t = \alpha^2 u_{xx}, && 0 < x < 1,\ 0 < t < \infty, \\
(\text{ГУ})\quad & u(0,t) = 0,\quad u_x(1,t) = 1, && 0 < t < \infty, \\
(\text{НУ})\quad & u(x,0) = \sin(\pi x), && 0 \leqslant x \leqslant 1.
\end{aligned}
\end{verbatim}

Чи можете ви графічно показати рішення цієї проблеми в різні моменти
часу? Чи буде це рішення, як правило, стаціонарним? Це очевидно?

\begin{enumerate}
\def\labelenumi{\arabic{enumi}.}
\setcounter{enumi}{2}
\tightlist
\item
  Яку фізичну інтерпретацію ви можете дати проблемі
\end{enumerate}

\begin{verbatim}
\begin{aligned}
(\text{УЧП})\quad & u_t = \alpha^2 u_{xx}, && 0 < x < 1,\ 0 < t < \infty, \\
(\text{ГУ})\quad & u_x(0,t) = 0,\quad u_x(1,t) = 0, && 0 < t < \infty, \\
(\text{НУ})\quad & u(x,0) = \sin(\pi x), && 0 \leqslant x \leqslant 1.
\end{aligned}
\end{verbatim}

Чи можете ви графічно показати рішення цієї проблеми в різні моменти
часу?

Що можна сказати про стабільну температуру?

\begin{enumerate}
\def\labelenumi{\arabic{enumi}.}
\setcounter{enumi}{3}
\tightlist
\item
  Припустимо, що бічна поверхня металевого стрижня не ізольована і має
  початкову температуру 20°C, але температура одного кінця стає 50°C.
  Потім стрижень промивають рідиною з температурою 30°C. Як виглядає
  змішана задача, що відповідає цьому випадку?
\end{enumerate}

\section{ВИВЕДЕННЯ РІВНЯННЯ
ТЕПЛОПРОВІДНОСТІ}\label{ux432ux438ux432ux435ux434ux435ux43dux43dux44f-ux440ux456ux432ux43dux44fux43dux43dux44f-ux442ux435ux43fux43bux43eux43fux440ux43eux432ux456ux434ux43dux43eux441ux442ux456}

МЕТА ЛЕКЦІЇ: Показати, як можна отримати одномірне рівняння
теплопровідності

\begin{verbatim}
u_t = \alpha^2 u_{xx} + f(x, t),
\end{verbatim}

спираючись лише на фундаментальний закон збереження кількості тепла.
Розглядається залежність швидкості теплопередачі від основних
термофізичних параметрів: теплопровідності, теплоємності та густини
матеріалу. Розглядаються різні модифікації базового рівняння
теплопровідності.

У кожній галузі знань існує набір базових принципів, правильність яких
очевидна. Усі інші твердження мають бути виведені з цих базових
принципів. Звісно, те, що очевидно одній людині, може здатися сумнівним
для іншої. Історія будь-якої науки полягає у постійному пошуку більш
фундаментальних принципів, що лежать в основі цієї науки. Основні
принципи мають бути настільки універсальними, щоб ніхто не сумнівався.

? --- Припущення A --- Припущення B --- Припущення C --- Припущення D
--- ? Рис. 4.1. Аксіоматичний метод.

Наприклад, одна людина може вважати, що всі факти певної науки можна
вивести з одного базового припущення, яке ми позначимо як B. На основі
припущення B вона або вона може довести теорему C, що призведе до
теореми D, а яка, у свою чергу, призведе до доведення багатьох інших
теорем (див. рисунок 4.1). Ось як виглядає прогрес у науці --- дослідник
може отримувати все більше нових результатів. Найвидатніші фізики,
хіміки та біологи або саме на цьому шляху.

З іншого боку, замість доведення нових теорем можна спробувати, якщо
можливо, знайти новий принцип (назвемо його A), який є більш
фундаментальним за B. Припущення A природно вважається більш
фундаментальним, ніж B.

Якщо B можна вивести з A. Отже, якщо можна знайти новий принцип A, тоді
межа наших знань зміщується до сфери більш фундаментальних ідей про
певну науку. У теорії теплопровідності основним принципом є закон
збереження енергії (теплова енергія). Усі інші твердження випливають із
цього базового принципу (рис. 4.2).

Закон збереження енергії

\begin{verbatim}
u_t = \alpha^2 u_{xx} + f(x,t)
\end{verbatim}

Інші властивості теплопровідності Рис. 4.2. Енергозбереження: наріжний
камінь теплопровідності.

Звісно, ми могли б забути цю лекцію і використати рівняння
теплопровідності як фундаментальне поняття (дехто може так подумати,
якщо це очевидно), але тоді ми обдуримо більш серйозних студентів, адже
закон збереження вважається найфундаментальнішим принципом у науці.
Науковці, моделюючи свої задачі, часто записують відношення, що
виражають закони збереження, а потім переписують їх у вигляді рівнянь у
похідних похідних.

Повернемося до головної мети нашої лекції --- вивести рівняння
теплопровідності з рівняння збереження енергії.

\section{Виведення рівняння
теплопровідності}\label{ux432ux438ux432ux435ux434ux435ux43dux43dux44f-ux440ux456ux432ux43dux44fux43dux43dux44f-ux442ux435ux43fux43bux43eux43fux440ux43eux432ux456ux434ux43dux43eux441ux442ux456-1}

Розглянемо однорідний стрижень довжини L, щодо якого зробимо такі
припущення:

\begin{enumerate}
\def\labelenumi{\arabic{enumi}.}
\tightlist
\item
  Стрижень виготовлений з одного однорідного провідного матеріалу.
\end{enumerate}

\pandocbounded{\includegraphics[keepaspectratio]{images/_page_29_Figure_10.jpeg}}

Фіг. 4.3. Тонкий тепловий стержень.

\begin{enumerate}
\def\labelenumi{\arabic{enumi}.}
\setcounter{enumi}{1}
\tightlist
\item
  Бікова поверхня стрижня термічно ізольована (тепло може поширюватися
  лише вздовж осі x).
\end{enumerate}

\texttt{3.\ Стрижень\ тонка,\ тобто\ температура\ всіх\ точок\ у\ кожному\ поперечному\ перерізі\ стрижня\ є\ сталою.\ Якщо\ розглядати\ частину\ вудки\ на\ сегменті\ \textquotesingle{}{[}x,\ x+\textbackslash{}Delta\ x{]}\textquotesingle{}\ і\ використовувати}

закон збереження кількості тепла, тоді можна написати:

\texttt{(4.1)\ Загальна\ зміна\ кількості\ тепла\ на\ сегменті\ \textquotesingle{}{[}x,\ x\ +\ \textbackslash{}Delta\ x{]}\ =\textquotesingle{}\ дорівнює\ загальній\ кількості\ тепла,\ що\ пройшло\ через\ межі\ +\ Нульова\ кількість\ тепла,\ утвореного\ всередині\ сегмента\ \textquotesingle{}{[}x,\ x\ +\ \textbackslash{}Delta\ x{]}\textquotesingle{}\ .}

\texttt{Оскільки\ загальна\ кількість\ тепла\ (у\ калоріях)\ всередині\ сегмента\ \textquotesingle{}\textbackslash{}\{x,\ x\ +\ \textbackslash{}Delta\ x\textbackslash{}\}\textquotesingle{}\ в\ будь-який\ момент\ часу\ t\ обчислюється\ формулою:}

\texttt{Загальна\ кількість\ тепла\ всередині\ сегмента\ \textquotesingle{}{[}x,\ x+\textbackslash{}Delta\ x{]}\ =\ \textbackslash{}int\textbackslash{}limits\_\{-\textbackslash{}infty\}\^{}\{x+\textbackslash{}Delta\ x\}\ coAu\textquotesingle{}\ (s,\ t)\ ds,}

де c --- питома теплоємність матеріалу (показує здатність матеріалу
накопичувати тепло),

Щільність матеріалу,

A --- це площа поперечного перерізу стрижня.

\textbf{va}con збереження енергії (4.1) можна надати у такій
математичній формі:

(4.2)

\begin{verbatim}
\frac{d}{dt} \int_{x}^{x+\Delta x} c\rho Au(s, t) ds = c\rho A \int_{x}^{x+\Delta x} u_{t}(s, t) ds = eA \left[ u_{x}(x+\Delta x, t) - u_{x}(x, t) \right] + A \int_{x+\Delta x}^{x+\Delta x} f(s, t) ds.
\end{verbatim}

\textbf{де} k --- теплопровідність матеріалу (вказує на здатність
матеріалу проводити тепло), \texttt{f(x,\ t)} --- об'ємна потужність
зовнішнього джерела тепла (калорії/см·с).

Задача полягає в тому, щоб записати рівняння (4.2) у формі, яка не
містить інтегралів. Щоб розв'язати цю задачу, нагадаємо читачеві теорему
про середнє значення з ходу інтегральної нечисельності.

\section{Теорема про
значення}\label{ux442ux435ux43eux440ux435ux43cux430-ux43fux440ux43e-ux437ux43dux430ux447ux435ux43dux43dux44f}

Якщо функція \texttt{f(x)} є неперервною на сегменті {[}a, b{]}, то
існує принаймні одна точка
\texttt{\textbackslash{}xi\ \textbackslash{}in\ {[}a,\ b{]}} така, що

\begin{verbatim}
\int_{a}^{b} f(x) dx = f(\xi) (b-a).
\end{verbatim}

Застосовуючи цей результат до рівняння (4.2), ми отримуємо наступний
результат

\begin{verbatim}
coAu_{t}(\xi, t) \Delta x = kA\{u_{x}(x + \Delta x, t) - u_{x}(x, t)\} + Af(\xi, t) \Delta x,
\end{verbatim}

\texttt{x\ \textless{}\ \textbackslash{}xi\ \textless{}\ x\ +\ \textbackslash{}Delta\ x,\ або\ \textquotesingle{}u\_t\textquotesingle{}(\textbackslash{}xi,\ t)\textquotesingle{}=\ \textbackslash{}frac\{k\}\{c\textbackslash{}rho\}\ \textbackslash{}left\textbackslash{}\{\ \textbackslash{}frac\{\textquotesingle{}u\_x\textquotesingle{}(x\ +\ \textbackslash{}Delta\ x,\ t)\ -\textquotesingle{}u\_x\textquotesingle{}(x,\ t)\}\{\textbackslash{}Delta\ x\}\ \textbackslash{}right\textbackslash{}\}\ +\ \textbackslash{}frac\{1\}\{c\textquotesingle{}\textbackslash{}rho\}\ f(\textbackslash{}xi,\ t).}

\begin{verbatim}
Спрямовуючи \Delta x до 0, отримуємо потрібне рівняння
\end{verbatim}

(4.3)

\begin{verbatim}
u_t(x, t) = \alpha^2 u_{xx}(x, t) + F(x, t),
\end{verbatim}

де

\begin{verbatim}
\alpha^{2}=\frac{k}{c\rho}
\end{verbatim}

\begin{verbatim}
— коефіцієнт теплової дифузії, F\left(x,\ t\right)=\frac{1}{c\rho}f\left(x,\ t\right) — густина джерел тепла.
\end{verbatim}

Ми могли б зупинитися на цьому. Однак ми також розглянемо випадок, коли
бічна поверхня стрижня не термічно ізольована. Припустимо, що кількість
теплового потоку через бічну поверхню стрижня в цьому випадку
пропорційна різниці між температурою стрижня \texttt{u(x,t)} і
температурою навколишнього середовища, яка залишається сталою і дорівнює
нулю. У цьому випадку закон збереження кількості тепла веде до рівняння

\begin{verbatim}
(4.4) u_t = \alpha^2 u_{xx} - \beta u + F(x, t),
\end{verbatim}

\begin{verbatim}
де \beta — коефіцієнт пропорційності потоку через бічну поверхню.
\end{verbatim}

\section{ЗАУВАЖЕННЯ}\label{ux437ux430ux443ux432ux430ux436ux435ux43dux43dux44f-3}

\begin{enumerate}
\def\labelenumi{\arabic{enumi}.}
\tightlist
\item
  Як зазначалося вище, константа k називається теплопровідністю речовини
  і чисельно дорівнює кількості тепла (у калоріях), що проходить за
  секунду через пластину товщиною 1 см і площею поперечного перерізу 1
  см³ при різниці температур на протилежних гранях пластини рівною 1°C.
  Числові значення теплопровідності для різних матеріалів можна знайти в
  довідниках з фізики або хімії. Типові значення теплопровідності
  варіюються від 1 (для міді) до значень, близьких до нуля для хороших
  теплоізоляторів. Якщо стрижень зроблений із однорідного матеріалу, то
  значення k не залежить від x. Для деяких матеріалів значення k
  залежить від температури, і тому рівняння теплопровідності в цьому
  випадку стає нелінійним:
\end{enumerate}

\begin{verbatim}
u_{t} = \frac{1}{c\rho} \frac{\partial}{\partial x} \{ k (u) | u_{x} \}.
\end{verbatim}

Однак найчастіше теплопровідність k дуже слабо впливає на температуру u,
і цю нелінійність можна ігнорувати.

\begin{enumerate}
\def\labelenumi{\arabic{enumi}.}
\setcounter{enumi}{1}
\item
  Значення c зазвичай називають питомою теплоємністю речовини. Він
  характеризує здатність речовини зберігати теплову енергію. Наприклад,
  запечена картопля має високу теплоємність, оскільки може зберігати
  велику кількість тепла на одиницю маси (саме тому вона зберігає високу
  температуру так довго). Чисельна теплоємність дорівнює кількості
  калорій тепла, які потрібно передати 1 граму речовини, щоб підвищити
  її температуру на 1°C. Для більшості розглянутих задач значення c
  можна вважати сталим, не записуючи ні x, ні i. Значення питомої
  теплоємності для конкретних матеріалів можна знайти у фізичних
  довідниках.
\item
  Одиниці вимірювання основних величин, включених у рівняння
\end{enumerate}

теплопровідність (у системі GHS) наступна:

i-температура (градуси Цельсія), \texttt{u\_t} --- швидкість зміни
температури (°C/s), \texttt{u\_x} --- нахил температурної кривої
(°C/см), \texttt{u\_\{xx\}} --- опуклість температурної кривої (°C/см²),

\begin{verbatim}
c = щільність тепла (калор/г \cdot °C)
\end{verbatim}

\begin{verbatim}
k — теплопровідність (кал/см \cdot при \cdot °C), \rho — щільність (г/<sup>см3</sup>), \alpha^2 — це теплова дифузія (см<sup>²</sup>/с).
\end{verbatim}

\begin{verbatim}
4. Слід зазначити, що теплова дифузія матеріалу \alpha^2 = \frac{k}{\epsilon \rho} прямо пропорційна теплопровідності матеріалу і обернено пропорційна густині \rho та питому теплоємності \epsilon ; Ми сподіваємося, що цей зв'язок відповідає інтуїтивним ідеям читача.
\end{verbatim}

\section{ЗАВДАННЯ}\label{ux437ux430ux432ux434ux430ux43dux43dux44f-3}

\begin{verbatim}
1. Підставити розмірність кожної з u,\ u_t,\ \dots величин у рівняння
\end{verbatim}

\begin{verbatim}
u_t = \alpha^2 u_{xx} - \beta u
\end{verbatim}

\begin{verbatim}
і переконайтеся, що всі терміни мають однаковий розмір {}^{\circ}\mathrm{C}/c .
\end{verbatim}

\begin{enumerate}
\def\labelenumi{\arabic{enumi}.}
\setcounter{enumi}{1}
\tightlist
\item
  Підставити розмірність кожної з величин у рівняння
\end{enumerate}

\begin{verbatim}
u_t = \alpha^2 u_{xx} - v u_x,
\end{verbatim}

\begin{verbatim}
де \upsilon вимірюється в одиницях швидкості, і переконайтеся, що всі терми мають однакову розмірність.
\end{verbatim}

3 Отримаємо рівняння теплопровідності

\begin{verbatim}
u_{t} = \frac{1}{c\rho} \frac{\partial}{\partial x} [h(x) u_{x}] + f(x, t)
\end{verbatim}

\begin{verbatim}
Для випадку, коли теплопровідність k\left( x\right) залежить від координати X.
\end{verbatim}

\begin{enumerate}
\def\labelenumi{\arabic{enumi}.}
\setcounter{enumi}{3}
\tightlist
\item
  Припустимо, що \texttt{u(x,\ t)} --- це значення концентрації речовини
  в потоку, що рухається зі швидкістю v. Припустимо, що концентрація
  змінюється внаслідок процесів дифузії та конвекції. Виведемо рівняння
\end{enumerate}

\begin{verbatim}
u_1 = \alpha^2 u_{xx} - v u_x,
\end{verbatim}

\begin{verbatim}
Виходячи з того, що в будь-який момент часу в певному [x, x + \Delta x] сегменті матерії не з'являється і не зникає.
\end{verbatim}

ПРИМІТКА. Закон збереження матерії виглядає так:

\begin{verbatim}
Зміни маси речовини на сегменті [x, x + \Delta x] = = Зміна внаслідок дифузії через межі + Зміна внаслідок переносу через межі.
\end{verbatim}

Лекція 5

\section{ЗМІННЕ
РОЗДІЛЕННЯ}\label{ux437ux43cux456ux43dux43dux435-ux440ux43eux437ux434ux456ux43bux435ux43dux43dux44f}

МЕТА ЛЕКЦІЇ: Познайомити читача з потужним методом розділення змінних і
показати, як цей метод можна використати для розв'язання добре відомої
проблеми дифузії. Оскільки цей метод створює труднощі для деяких учнів
(головним чином через громіздкі алгебранічні перетворення), ми
супроводжуємо презентацію інтуїтивними міркуваннями.

Основна ідея методу полягає в тому, щоб розкласти початкову умову на
найпростіші компоненти, знайти відгук системи на кожен найпростіший
компонент і підсумувати всі відповіді. Таким чином можна знайти
відповідь на довільну початкову умову.

Існуючий метод поділу змінних певною мірою зменшує можливість такої
інтерпретації методу, але все ж не виключає його.

Метод розділення змінних є одним із найвідоміших методів розв'язання
зсувних задач і використовується, якщо:

\begin{itemize}
\tightlist
\item
  \begin{enumerate}
  \def\labelenumi{\arabic{enumi}.}
  \tightlist
  \item
    Рівняння є лінійним і однорідним (не обов'язково з постійними
    коефіцієнтами).
  \end{enumerate}
\item
  \begin{enumerate}
  \def\labelenumi{\arabic{enumi}.}
  \setcounter{enumi}{1}
  \tightlist
  \item
    Граничні умови визначаються як
  \end{enumerate}
\end{itemize}

\begin{verbatim}
\alpha u_x(0, t) + \beta u(0, t) = 0, \qquad \gamma u_x(1, t) + \delta u(1, t) = 0.
\end{verbatim}

\begin{verbatim}
де \alpha, \beta, \gamma і \delta — константи (граничні умови, задані в такій формі, називаються лінійними однорідними граничними умовами).
\end{verbatim}

Метод був створений за часів Фур'є (зазвичай відомий як метод Фур'є) і,
здається, є найпопулярнішим (де це застосовно).

\pandocbounded{\includegraphics[keepaspectratio]{images/_page_34_Figure_8.jpeg}}

Рис. 5.1. Схематичне представлення задачі дифузії.

Замість того, щоб вивчати метод у загальному випадку, давайте спочатку
розглянемо конкретну задачу (загальний випадок обговоримо пізніше).
Розглянемо змішану задачу типу дифузії: щоб знайти розв'язок

(PPP)

\begin{verbatim}
u_t = \alpha^2 u_{xx}
, 0 < x < 1 , 0 < t < \infty ,
\end{verbatim}

що задовольняє граничні умови

(GI)

\begin{verbatim}
\begin{cases} u(0, t) = 0, & 0 < t < \infty, \\ u(1, t) = 0, & 0 < t < \infty, \end{cases}
\end{verbatim}

та початкова умова

(ГУ)

\begin{verbatim}
u(x, 0) = \varphi(x), 0 \le x \le 1.
\end{verbatim}

\begin{verbatim}
Перш ніж почати розділяти змінні, давайте дамо фізичну інтерпретацію нашої проблеми. Отже, існує стрижень скінченної довжини, кінці якого підтримуються при сталій температурі, рівній нулю (насправді кінці можна підтримувати при значно вищій температурі, значення якої вважається початковою точкою). Додаткові дані про завдання подаються як початкова умова. Наша мета — знайти розподіл температури u\left(x,t\right) в наступний час.
\end{verbatim}

\section{Загальні принципи методу розділення
змінних}\label{ux437ux430ux433ux430ux43bux44cux43dux456-ux43fux440ux438ux43dux446ux438ux43fux438-ux43cux435ux442ux43eux434ux443-ux440ux43eux437ux434ux456ux43bux435ux43dux43dux44f-ux437ux43cux456ux43dux43dux438ux445}

Для найпростішого диференціального рівняння в похідних похідних
розділення змінних є пошуком розв'язків у вигляді

\begin{verbatim}
u(x, t) = X(x)T(t),
\end{verbatim}

де \texttt{X(x)} --- це функція, що залежить лише від змінної x, а
\texttt{T(t)} --- функція, що залежить лише від t. Це розв'язок у
певному сенсі є найпростішим, оскільки температура \texttt{u(x,\ t)},
представлена у такій формі, зберігатиме «форму» профілю у різний час
часу t (див. рисунок 5.2).

\pandocbounded{\includegraphics[keepaspectratio]{images/_page_35_Figure_6.jpeg}}

РИСУНОК 5.2 Графік функції X(x)

\begin{verbatim}
T(t) у різних точках часу t.

Загальна ідея полягає в тому, щоб знайти нескінченну кількість таких часткових диференціальних розв'язків (які задовольняють граничні умови
\end{verbatim}

\texttt{u(0,\ t)\ =\ 0,\ \textbackslash{}qquad\ u(1,\ t)\ =\ 0.}

\begin{verbatim}
Для цього підставляємо розв'язки (5.1) у ці граничні умови. У результаті отримуємо
\end{verbatim}

\texttt{u(0,\ t)\ =\ B\ e\^{}\{-\textbackslash{}lambda\^{}2\ \textbackslash{}alpha\^{}2\ t\}\ =\ 0\ \textbackslash{}Rightarrow\ B\ =\ 0,}
\texttt{u(1,\ t)\ =\ A\ e\^{}\{-\textbackslash{}lambda\^{}2\ \textbackslash{}alpha\^{}2\ t\}\textbackslash{}sin\ \textbackslash{}lambda\ =\ 0\ \textbackslash{}Rightarrow\ \textbackslash{}sin\ \textbackslash{}lambda\ =\ 0.}

\begin{verbatim}
Отже,
\end{verbatim}

\texttt{\textbackslash{}lambda\_n\ =\ n\textbackslash{}pi,\ \textbackslash{}qquad\ n\ =\ 1,\ 2,\ \textbackslash{}ldots}

\begin{verbatim}
Зверніть увагу, що друга гранична умова могла б виконуватися і при A=0 , але тоді розв'язок (5.1) був би тривіально нульовим.

![](images/_page_38_Figure_3.jpeg)

Фіг. 5.3. Фундаментальне рішення n_n(x, t) = A_n e^{-(\pi n \alpha)^{n/t}} \sin{(\pi n x)} .

Тепер, коли ми зробили другий крок, маємо нескінченний набір ознак

(5.2)
\end{verbatim}

\texttt{u\_n(x,\ t)\ =\ A\_n\ e\^{}\{-(n\textbackslash{}pi\textbackslash{}alpha)\^{}2\ t\}\ \textbackslash{}sin(n\textbackslash{}pi\ x),\ \textbackslash{}quad\ n\ =\ 1,\ 2\ \textbackslash{}dots}

\begin{verbatim}
кожна з яких задовольняє рівняння з частковими похідними та граничними умовами). Отримані функції — це будівельні блоки, з яких ми будемо створювати рішення нашої проблеми. Це рішення буде певною сумою

<sup>2)</sup> Зверніть увагу, що функції u_n і u_{-n} відрізняються лише знаком.

цих найпростіших функцій. Конкретний тип суми залежить від початкового стану. Фіг. Рисунок 5.3 показує графіки кількох таких фундаментальних функцій

КРОК 3. (Пошук розв'язку, який задовольняє рівняння, межу та початкові умови.)

Останній (і, ймовірно, найцікавіший з математичної точки зору) крок — знайти таку суму фундаментальних розв'язків
\end{verbatim}

\texttt{u(x,\ t)\ =\ \textbackslash{}sum\_\{n=1\}\^{}\{\textbackslash{}infty\}\ A\_n\ e\^{}\{-(n\textbackslash{}pi\textbackslash{}alpha)\^{}n\ t\}\ \textbackslash{}sin(n\textbackslash{}pi\ x),}

\begin{verbatim}
тобто при виборі таких коефіцієнтів A_n , що функція задовольнятиме початкову умову
\end{verbatim}

\texttt{u\textbackslash{}left(\ x,\textbackslash{}\ 0\textbackslash{}right)\ =\textbackslash{}varphi\textbackslash{}left(\ x\textbackslash{}right)\ .}

\begin{verbatim}
Підставляючи суму в початковому стані, отримуємо

(5.3)
\end{verbatim}

\texttt{q(x)\ =\ \textbackslash{}sum\_\{n=1\}\^{}\{\textbackslash{}infty\}\ A\_n\ \textbackslash{}sin(n\textbackslash{}pi\ x).}

\begin{verbatim}
Це рівняння підводить нас до цікавого питання: чи можна розкласти початкову температуру \phi(x) на ряд елементарних функцій вигляду
\end{verbatim}

\texttt{A\_1\ \textbackslash{}sin(\textbackslash{}pi\ x)\ +\ A\_2\ \textbackslash{}sin(2\textbackslash{}pi\ x)\ +\ A\_3\ \textbackslash{}sin(3\textbackslash{}pi\ x)\ +\ \textbackslash{}dots}

\begin{verbatim}
?

Позитивну відповідь на це питання дав французький математик Жозеф Фур'є. Виявилося, що для достатньо хороших функцій така декомпозиція можлива ^1 ). Тоді виникає нове питання: знайти коефіцієнти розкладу A_n ?

Насправді це дуже просто. Для цього використаємо властивість системи функцій
\end{verbatim}

\texttt{\textbackslash{}\{\textbackslash{}sin(n\textbackslash{}pi\ x)\textbackslash{}\}\_\{n=1\}\^{}\{\textbackslash{}infty\}.}

\begin{verbatim}
відомий як ортогональність.

У нашому випадку (див. завдання 2) ці функції задовольняють умови
\end{verbatim}

\texttt{\textbackslash{}int\_\{0\}\^{}\{1\}\ \textbackslash{}sin(m\textbackslash{}pi\ x)\ \textbackslash{}sin(n\textbackslash{}pi\ x)\ dx\ =\ \textbackslash{}begin\{cases\}\ 0,\ \&\ m\ \textbackslash{}neq\ n,\ \textbackslash{}\textbackslash{}\ 1/2,\ \&\ m\ =\ n.\ \textbackslash{}end\{cases\}}

\begin{verbatim}
Умови розкладу функцій у рядах Фур'є наведені в лекції 11. Для глибшого ознайомлення з цим питанням можна звернутися до численних освітніх літератур (див., наприклад: Г. П. Толстов, Серія Фур'є. — Москва: Наука, 1980). — Ед. Ед.

Ця властивість чітко показана на рис. 5.4, яка показує графіки деяких функцій із цієї системи.

Отже, ми хочемо знайти коефіцієнти у розкладі
\end{verbatim}

\texttt{\textbackslash{}varphi(x)\ =\ A\_1\ \textbackslash{}sin(\textbackslash{}pi\ x)\ +\ A\_2\ \textbackslash{}sin(2\textbackslash{}pi\ x)\ +\ A\_3\ \textbackslash{}sin(3\textbackslash{}pi\ x)\ +\ \textbackslash{}dots}

\begin{verbatim}
![](images/_page_40_Figure_5.jpeg)

РИСУНОК 5.4. Послідовність ортогональних функцій

Помножте обидві сторони цього співвідношення на \sin{(m\pi x)} (m — довільне ціле число) і інтегруйте від нуля до одиниці. У результаті ми отримуємо
\end{verbatim}

\texttt{\textbackslash{}int\_\{0\}\^{}\{1\}\ \textbackslash{}varphi(x)\ \textbackslash{}sin(m\textbackslash{}pi\ x)\ \textbackslash{},\ dx\ =\ A\_\{m\}\ \textbackslash{}int\_\{0\}\^{}\{1\}\ \textbackslash{}sin\^{}\{2\}(m\textbackslash{}pi\ x)\ \textbackslash{},\ dx\ =\ \textbackslash{}frac\{A\_\{m\}\}\{2\}}

\begin{verbatim}
(усі інші члени зрівнялися з нулем через ортогональність). Розв'язуючи рівняння відносно A_m , отримуємо
\end{verbatim}

\texttt{A\_m\ =\ 2\ \textbackslash{}int\_0\^{}1\ \textbackslash{}varphi(x)\ \textbackslash{}sin(m\textbackslash{}pi\ x)\ dx.}

\begin{verbatim}
Отже, отримано, що розв'язок записується у вигляді

(5.4)
\end{verbatim}

\texttt{u(x,\ t)\ =\ \textbackslash{}sum\_\{n=1\}\^{}\{\textbackslash{}infty\}\ A\_n\ e\^{}\{-(n\textbackslash{}pi\textbackslash{}alpha)\^{}2\ t\}\ \textbackslash{}sin(\textbackslash{}pi\ n\ x),}

\begin{verbatim}
де коефіцієнти A_n визначаються формулами
\end{verbatim}

\texttt{A\_n\ =\ 2\textbackslash{}int\_0\^{}1\ \textbackslash{}varphi(x)\ \textbackslash{}sin(\textbackslash{}pi\ nx)\ dx}

\begin{verbatim}
Ви можете переконатися, що отримане рішення задовольняє всі умови початкової задачі. На цьому завершується третій крок.

Багато студентів розчаровані, коли бачать, наскільки складним виглядає рішення, і майже не намагаються дивитися на нього знову (виглядає дуже погано). Однак це рішення не здаватиметься таким складним, якщо ви витратите час на його аналіз. Насправді, складніша форма розв'язку робить його більш інформативним. Нижче наведено кілька нотаток, які допоможуть вам інтерпретувати це рішення.

# ЗАУВАЖЕННЯ

Зверніть увагу, що єдина різниця між синусоїдальним розкладом функції φ(x) і розв'язком (5.4) — це наявність часового фактора e<sup>-(nπα)<sup>3</sup>1</sup> у кожному члені ряду. Отже, якщо початкова умова має просту декомпозицію, наприклад, вигляду
\end{verbatim}

\texttt{\textbackslash{}varphi(x)\ =\ \textbackslash{}sin(\textbackslash{}pi\ x)\ +\ \textbackslash{}frac\{1\}\{2\}\textbackslash{}sin(3\textbackslash{}pi\ x),}

\begin{verbatim}
тоді рішення можна негайно записати
\end{verbatim}

\texttt{u(x,\ t)\ =\ e\^{}\{-(\textbackslash{}pi\textbackslash{}alpha)\^{}2\ t\}\ \textbackslash{}sin(\textbackslash{}pi\ x)\ +\ \textbackslash{}frac\{1\}\{2\}\ e\^{}\{-(3\textbackslash{}pi\textbackslash{}alpha)\^{}2\ t\}\ \textbackslash{}sin(3\textbackslash{}pi\ x).}

\begin{verbatim}
У цьому випадку очевидно, що якщо ми розкладемо \phi(x) на ряд Фур'є вздовж синусів, то отримуємо
\end{verbatim}

A\_1 = 1,

\begin{verbatim}
A_2 = 0, A_3 = 1/2, A_4 = A_5 = \dots = 0. - 2. Розв'язок (5.4) можна інтерпретувати так: ми представляємо початкову температуру \varphi(x) як суму найпростіших функцій виду A_n \sin(\pi nx) ; кожна така функція породжує відповідь A_n e^{-(\pi n\alpha)^2 t} \sin(\pi nx) . Сумуючи всі такі відповіді, отримуємо розв'язок, який відповідає початковій умові u(x, 0) = \varphi(x) .
- 3. Кожен член у декомпозиції
\end{verbatim}

\texttt{u(x,\ t)\ =\ A\_1\ e\^{}\{-(\textbackslash{}pi\textbackslash{}alpha)\^{}2\ t\}\ \textbackslash{}sin(\textbackslash{}pi\ x)\ +\ A\_2\ e\^{}\{-(2\textbackslash{}pi\textbackslash{}alpha)\^{}2\ t\}\ \textbackslash{}sin(2\textbackslash{}pi\ x)\ +\ \textbackslash{}cdots}

\begin{verbatim}
є функцією від x і t. Зверніть увагу, що внесок членів з великими числами в t>0 дуже малий через множник e^{-(n\pi\alpha)^2 t} . Відповідно, після закінчення досить великого Р и с

![](images/_page_42_Figure_2.jpeg)

. 5.5. У задачах дифузії вищі порядки терміни швидко затухають; \phi(x) — початкова температура. Слід зазначити, що по-перше, у розчині зникають високі частоти — по-перше, зникли найбільші коливання міді.

Повний розв'язок приблизно збігається з першим членом, який є півхвилі синусоїди, що з часом згасає (рис. 1). 5.5.)

# ЗАВДАННЯ

1. Показати, що функції
\end{verbatim}

\texttt{u(x,\ t)\ =\ e\^{}\{-\textbackslash{}lambda\^{}2\ \textbackslash{}alpha\textquotesingle{}\ t\}\ \textbackslash{}left{[}\ A\ \textbackslash{}sin(\textbackslash{}lambda\ x)\ +\ B\ \textbackslash{}cos(\textbackslash{}lambda\ x)\ \textbackslash{}right{]}}

\begin{verbatim}
задовольняють рівняння u_l = \alpha^2 u_{xx} для довільних значень A, B і \lambda .

2. Покажи це
\end{verbatim}

\texttt{\textbackslash{}int\_0\^{}1\ \textbackslash{}sin(m\textbackslash{}pi\ x)\ \textbackslash{}sin(n\textbackslash{}pi\ x)\ dx\ =\ \textbackslash{}begin\{cases\}\ 0,\ \&\ m\ \textbackslash{}neq\ n,\ \textbackslash{}\textbackslash{}\ 1/2,\ \&\ m\ =\ n.\ \textbackslash{}end\{cases\}}

\begin{verbatim}
ПРИМІТКА. Використання ідентичності
\end{verbatim}

\texttt{\textbackslash{}sin(mx)\textbackslash{}sin(nx)\ =\ \textbackslash{}frac\{1\}\{2\}{[}\textbackslash{}cos(m-n)x\ -\ \textbackslash{}cos(m+n)x{]}.}

\begin{verbatim}
3. Знайти розклад у рядах Фур'є синусами функції \varphi(x) = 1 на сегменті [0, 1]. Побудуйте перші три терміни розкладу.

Використовуючи результати завдання 3, знайдіть розв'язок наступної змішаної задачі:

(PPP)
\end{verbatim}

u\_t = u\_\{xx\}

\begin{verbatim}
0 < x < 1 , 
(GI)
\end{verbatim}

\texttt{\textbackslash{}begin\{cases\}\ u(0,\ t)\ =\ 0,\ \textbackslash{}\textbackslash{}\ u(1,\ t)\ =\ 0,\ \textbackslash{}end\{cases\}}

\begin{verbatim}
0 < t < \infty , 
(НУ) u(x, 0) = 1, \quad 0 \le x \le 1 .

(Зверніть увагу, що ця проблема фізично не має значення, оскільки припускається, що температура на кінцях стрижня миттєво знижується з одного до нуля \phi(x) .

5. Знайти розв'язок проблеми 4, якщо початкова умова встановлена у вигляді
\end{verbatim}

\texttt{u(x,\ 0)\ =\ \textbackslash{}sin(2\textbackslash{}pi\ x)\ +\ \textbackslash{}frac\{1\}\{3\}\textbackslash{}sin(4\textbackslash{}pi\ x)\ +\ \textbackslash{}frac\{1\}\{5\}\textbackslash{}sin(6\textbackslash{}pi\ x).}

\begin{verbatim}
6. Розв'яжіть задачу 4, якщо початкова умова u\left( x,\ 0\right) =x-x^{2},\quad 0< x<1. Лекція 6

# ПЕРЕТВОРЕННЯ НЕОДНОРІДНИХ ГРАНИЧНИХ УМОВ У ОДНОРІДНІ ГРАНИЧНІ УМОВИ

МЕТА ЛЕКЦІЇ: Показати, як виникає змішана задача

~~
\begin{aligned}
(\text{УЧП})\quad & u_t - a^2 u_{xx} = f(x,t), \\
(\text{ГУ})\quad & \alpha_1 u_x(0,t) + \beta_1 u(0,t) = g_1(t), \\
& \alpha_2 u_x(L,t) + \beta_2 u(L,t) = g_2(t), \\
(\text{НУ})\quad & u(x,0) = \varphi(x).
\end{aligned}
~~

можна звести до задачі вигляду
\end{verbatim}

\texttt{U\_\{t\}\ -\ \textbackslash{}alpha\^{}\{2\}U\_\{xx\}\ =\ F(x,\ t),}
\texttt{\textbackslash{}alpha\_\{1\}U\_\{x\}(0,\ t)\ +\ \textbackslash{}beta\_\{1\}U(0,\ t)\ =\ 0,}
\texttt{\textbackslash{}alpha\_\{2\}U\_\{x\}(L,\ t)\ +\ \textbackslash{}beta\_\{2\}U(L,\ t)\ =\ 0,}
\texttt{U(x,\ 0)\ =\ \textbackslash{}varphi(x)}

\begin{verbatim}
з нульовими граничними умовами. Цей новий виклик можна подолати

- 1. Шляхом ділення змінних (якщо нова задача виявиться однорідною, тобто F(x, t) = 0).

~~
- 2. Методом інтегральних перетворень і розкладу власними функціями, якщо F(x, t) \neq 0 .
~~

Метод розділення змінних, про який йшлося в попередній лекції, є досить потужним, а розв'язки, отримані за його допомогою, подаються у зручній формі. Однак читач повинен розуміти, що цей метод не застосовний до всіх задач. Щоб метод розділення змінних був застосовним, граничні умови мають бути лінійними, однорідними, тобто представленими як

(6.1)
\end{verbatim}

\texttt{\textbackslash{}begin\{aligned\}\ \textbackslash{}alpha\_1\ u\_x(0,\ t)\ +\ \textbackslash{}beta\_1\ u(0,\ t)\ \&=\ 0,\ \textbackslash{}\textbackslash{}\ \textbackslash{}alpha\_2\ u\_x(L,\ t)\ +\ \textbackslash{}beta\_2\ u(L,\ t)\ \&=\ 0.\ \textbackslash{}end\{aligned\}}

\begin{verbatim}
Мета цієї лекції — показати, як виникає задача з неоднорідними граничними умовами форми

(6.2)
\end{verbatim}

\texttt{\textbackslash{}begin\{aligned\}\ \&\ (\textbackslash{}mathsf\{Y\}\textbackslash{}mathsf\{Ch\}\textbackslash{}Pi)\ \&\ u\_t\ =\ \textbackslash{}alpha\^{}2\ u\_\{xx\},\ \textbackslash{}\textbackslash{}\ \&\ (\textbackslash{}mathsf\{\textbackslash{}Gamma\}\textbackslash{}mathsf\{Y\})\ і\ \textbackslash{}begin\{case\}\ \textbackslash{}alpha\_1\ u\_x\ \textbackslash{},\ (0,\ \textbackslash{},\ t)\ +\ \textbackslash{}beta\_1\ u\ \textbackslash{},\ (0,\ \textbackslash{},\ t)\ =\ g\_1\ \textbackslash{},\ (t)\ \textbackslash{}\textbackslash{}\ \textbackslash{}alpha\_2\ u\_x\ \textbackslash{},\ (L,\ \textbackslash{},\ t)\ +\ \textbackslash{}beta\_2\ u\ \textbackslash{},\ (L,\ \textbackslash{},\ t)\ =\ g\_2\ \textbackslash{},\ (t)\ \textbackslash{}end\{case\}\ \textbackslash{}text\{\ (гетерогенний\ \}\ \textbackslash{}mathsf\{\textbackslash{}Gamma\}\textbackslash{}mathsf\{Y\}),\ \textbackslash{}\textbackslash{}\ \&\ (\textbackslash{}mathsf\{H\}\textbackslash{}mathsf\{Y\})\ і\ u\ \textbackslash{},\ \ (x,\ \textbackslash{},\ 0)\ =\ \textbackslash{}varphi\ \textbackslash{},\ (x)\ \textbackslash{}end\{aligned\}}

\begin{verbatim}
можна розв'язати, зведивши її до задачі з однорідними граничними умовами.

Нове завдання можна розв'язати іншими методами

(наприклад, шляхом розкладу за власними функціями).

Почнемо з надзвичайно простого прикладу перетворення гетерогенних граничних умов на однорідні.

# Перетворення неоднорідних граничних умов у однорідні граничні умови

Розглянемо проблему поширення тепла в термоізольованому стрижні, кінці якого підтримуються при сталих температурах k_1 і k_2, тобто
\end{verbatim}

\texttt{\textbackslash{}begin\{aligned\}}
\texttt{(\textbackslash{}text\{УЧП\})\textbackslash{}quad\ \&\ u\_t\ =\ \textbackslash{}alpha\^{}2\ u\_\{xx\},\ \&\&\ 0\ \textless{}\ x\ \textless{}\ L,\textbackslash{}\ 0\ \textless{}\ t\ \textless{}\ \textbackslash{}infty,\ \textbackslash{}\textbackslash{}}
\texttt{(6.3)\textbackslash{}\ (\textbackslash{}text\{ГУ\})\textbackslash{}quad\ \&\ u(0,t)\ =\ k\_1,\textbackslash{}quad\ u(L,t)\ =\ k\_2,\ \&\&\ 0\ \textless{}\ t\ \textless{}\ \textbackslash{}infty,\ \textbackslash{}\textbackslash{}}
\texttt{(\textbackslash{}text\{НУ\})\textbackslash{}quad\ \&\ u(x,0)\ =\ \textbackslash{}varphi(x),\ \&\&\ 0\ \textbackslash{}leqslant\ x\ \textbackslash{}leqslant\ L.}
\texttt{\textbackslash{}end\{aligned\}}

\begin{verbatim}
Складність цієї задачі полягає в тому, що, оскільки граничні умови в ній є гетерогенними, ми не можемо розв'язати її шляхом розділення змінних. Однак очевидно, що в t \to \infty розв'язок нашої задачі схиляється до стаціонарного розв'язку, який лінійно (вздовж x) змінюється від температури k_1 до температури k_2 (рис. 6.1). Іншими словами, розумно припустити, що температуру в нашій задачі можна представити як суму
~~

# з двох термів
\end{verbatim}

\texttt{u\textbackslash{}left(x,\textbackslash{},t\textbackslash{}right)=}

\begin{verbatim}
~~
стаціонарна + тимчасова = \uparrow \uparrow Частина розв'язку, що залежить від HY і має тенденцію до нуля з помноженням на =\left[k_1+\frac{x}{L}\left(k_2-k_1\right)\right]+U\left(x,\,t\right). У цьому випадку наше завдання — знайти температуру переходу U\left( x,\,t\right) . Заміна
~~
\end{verbatim}

\texttt{u(x,\ t)\ =\ \textbackslash{}left{[}k\_1\ +\ \textbackslash{}frac\{x\}\{L\}(k\_2\ -\ k\_1)\textbackslash{}right{]}\ +\ U(x,\ t)}

\begin{verbatim}
~~
У початковій задачі (6.3) ми отримуємо нову задачу відносно невідомої функції U\left(x,\ t\right) . Ви можете вирішити проблему Р и с
~~

![](images/_page_45_Figure_7.jpeg)

~~
. 6.1. Розв'язання задачі (6.3) у різний час часу a- початкова температура u\left(x,\ 0\right)=\phi\left(x\right); b - стаціонарна температура u\left(x,\ \infty\right)=k_{1}+\frac{x}{L}\left(k_{2}-k_{1}\right); b - перехідна температура.
~~

~~
Відносно нова невідома функція U\left(x,\,t\right) і додає її до стаціонарного розв'язку, що дає бажану функцію u\left(x,\,t\right) . Виконавши ці прості перетворення з (6.3),
~~

Ми отримаємо

(6.4)

~~
\begin{aligned}
U_t &= \alpha^2 U_{xx}, && 0 < x < L, \\
U(0,t) &= 0,\quad U(L,t) = 0, && 0 < t < \infty, \\
U(x,0) &= \varphi(x) - \left[k_1 + \frac{x}{L}(k_2 - k_1)\right] = \overline{\varphi}(x).
\end{aligned}
~~

~~
де \overline{\phi}(x) — нова, але відома початкова умова.
~~

Ця задача стосується не лише однорідного рівняння, а, на щастя, однорідних граничних умов, що дозволяє розв'язувати її методом розділення змінних. Насправді, читач, ймовірно, пам'ятає рішення:

(6.5)
\end{verbatim}

\texttt{U(x,\ t)\ =\ \textbackslash{}sum\_\{n=1\}\^{}\{\textbackslash{}infty\}\ a\_n\ e\^{}\{-(n\textbackslash{}pi\textbackslash{}alpha)\^{}4\ t\}\ \textbackslash{}sin(n\textbackslash{}pi\ x/L),}

\begin{verbatim}
рей
\end{verbatim}

\texttt{a\_n\ =\ \textbackslash{}frac\{2\}\{L\}\ \textbackslash{}int\_\{-\textbackslash{}infty\}\^{}\{L\}\ \textbackslash{}overline\{\textbackslash{}phi\}(\textbackslash{}xi)\ \textbackslash{}sin(\textbackslash{}pi\ n\textbackslash{}xi/L)\ d\textbackslash{}xi.}

\begin{verbatim}
Це все для сердечників із фіксованою температурою на кінцях. А як щодо більш реалістичних GU з правою стороною, залежною від часу? Основні ідеї тут ті ж, що й у попередній задачі, але вони набагато складніші.

# Перетворення граничних умов, залежних від часу, у нульові

Розглянемо типову задачу

(PPP)
\end{verbatim}

\texttt{u\_t\ =\ \textbackslash{}alpha\^{}2\ u\_\{xx\}}

\begin{verbatim}
~~
, 0 < x < L , 0 < t < \infty , 0 < t < \infty , 0 < t < \infty , 0 < t < \infty , 0 < t < \infty , 0 < t < \infty , 0 < t < \infty , 0 < t < \infty , 0 < t < \infty , 0 < t < \infty , 0 < t < \infty , 0 < t < \infty , 0 < t < \infty , 0 < t < \infty , 0 < t < \infty , 0 < t < \infty , 0 < t < \infty , 0 < t < \infty , 0 < t < \infty , 0 < t < \infty , 0 < t < \infty , 0 < t < \infty , 0 < t < \infty , 0 < t < \infty , 0 < t < \infty , 0 < t < \infty , , 0 < t < \infty , 0 < t < \infty , 0 < t < \infty , 0 < t < \infty , 0 < t < \infty , 0 < t < \infty , 0 < t < \infty , 0 < t < \infty , 0 < t < \infty , 0 < t < \infty , 0 < t < \infty , 0 < t < \infty , 0 < t < \infty , 0 < t < \infty , 0 < t < \infty , 0 < t < \infty , 0 < t < \infty , 0 < t < \infty , 0 < t < \infty , 0 < t < \infty , 0 < t < \infty , 0 < t < \infty , 0 < t < \infty , 0 < t < \infty , 0 < t < \infty , 0 < t < \infty , 0 < t < \infty , 0 < t < \infty , 0 < t < \infty , 0 < t < \infty , 0 < t < \infty , 0 < t < \infty , 0 < t < \infty , 0 < t < \infty , 0 < t < \infty , 0 < t < \infty , 0 < t < \infty , 0 < t < \infty , 0 < t < \infty , 0 < t < \infty , 0 < t < \infty , 0 < t < \infty , 0 < t < \infty , 0 < t < \infty , 0 < t < \infty , 0 < t < \infty , 0 < t < \infty , 0 < t < \infty , 0 < t < \infty , 0 < t < \infty , 0 < t < \infty , 0 < t < \infty , 0 < t < \infty , 0 < t < \infty , 0 < t < \infty , 0 < t < \infty , 0 < t < \infty , 0 < t < \infty , 0 < t < \infty , 0 < t < \infty , 0 < t < \infty , 0 < t < \infty , 0 < t < \infty , 0 < t < \infty , 0 < t < \infty , 0 < t < \infty , 0 < t < \infty , 0 < t < \infty , 0 < t < \infty , 0 < t < \infty , 0 < t < \infty , 0 < t < \infty , 0 < t < \infty , 0 < t < \infty , 0 < t < \infty , 0 < t < \infty , 0 < t < \infty , 0 < t < \infty , 0 < t < \infty , 0 < t < \infty , 0 < t < \infty , 0 < t < \infty , 0 < t < \infty , 0 < t < \infty , 0 < t < \infty , 0 < t < \infty , 0 < t < \infty , 0 < t < \infty , 0 < t < \infty , 0 < t < \infty , 0 < t < \infty , 0 < t < \infty , 0 < t < \infty , 0 < t < \infty , 0 < t < \infty , 0 < t < \infty , 0 < t < \infty , 0 < t < \infty , 0 < t < \infty , 0 < t < \infty , 0 < t < \infty , 0 < t < \infty , 0 < t < \infty , 0 < t < \infty , 0 < t < \infty , 0 < t < \infty , 0 < t < \infty , 0 < t < \infty , 0 < t < \infty , 0 < t < \infty , , 0 < t < \infty , 0 < t < \infty , 0 < t < \infty , 0 < t < \infty , 0 < t < \infty , 0 < t < \infty , 0 < t < \infty , 0 < t < \infty , 0 < t < \infty , 0 < t < \infty , 0 < t < \infty , 0 < t < \infty , 0 < t < \infty , 0 < t < \infty , 0 < t < \infty , 0 < t < \infty , 0 < t < \infty , 0 < t < \infty , 0 < Щоб перетворити ці граничні умови на нульові умови, ми (після деяких спроб і помилок) зупинилися на такій формі розв'язку:
~~

(6.7)
\end{verbatim}

u(x, t) = A(t){[}1-x/L{]} + B(t){[}x/L{]} + U(x, t),

\begin{verbatim}
~~
де функції A\left(t\right) і B\left(t\right) обираються так, що «квазі-стаціонарна» частина розв'язку (6.7)
~~

(6.8)
\end{verbatim}

S(x, t) = A(t){[}1 - x/L{]} + B(t){[}x/L{]}

\begin{verbatim}
~~
задовольняло граничні умови початкової задачі. У цьому випадку функція U\left(x,\ t\right) задовольнятиме однорідні граничні умови.
~~

Підставляючи функцію S(x, t) у граничні умови
\end{verbatim}

S(0, t) = g\_1(t),

\begin{verbatim}
S_x(L, t) + hS(L, t) = g_2(t) призводить до двох рівнянь, з яких можна визначити A(t) і B(t). У результаті ми отримуємо

(6.9)
\end{verbatim}

\texttt{A(t)\ =\ g\_1(t),\ B(t)\ =\ \textbackslash{}frac\{g\_1(t)\ +\ Lg\_2(t)\}\{1\ +\ Lh\}.}

\begin{verbatim}
Отже.
\end{verbatim}

\texttt{u(x,\ t)\ =\ g\_1(t)\ {[}1\ -\ x/L{]}\ +\ \textbackslash{}frac\{g\_1(t)\ +\ Lg\_2(t)\}\{1\ +\ Lh\}\ {[}x/L{]}\ +\ U(x,\ t).}

\begin{verbatim}
Якщо ми підставимо ці вирази для you(x, t) у початкову задачу (це просять читача), отримаємо нову задачу для невідомої функції U(x, t):

~~
\begin{array}{llllllllllllllllllllllllllllllllllll
~~

Тепер ми стикаємося з новою проблемою однорідних граничних умов (на жаль, рівняння стало гетерогенним). Ми не можемо розв'язати цю задачу методом розділення змінних, але якщо читач погоджується чекати, то за кілька лекцій він ознайомиться з розв'язком цієї задачі методом цілісних перетворень і розкладання за власними функціями.

# ЗАУВАЖЕННЯ .

1. Як зазначалося раніше, мета цієї лекції — навчити читача, як перетворювати задачі з неунірольними граничними умовами на задачі з нульовими граничними умовами. Після цього може статися, що нове диференціальне рівняння в похідних похідних стане однорідним. Якщо це станеться, ми можемо вважати себе щасливцями, оскільки в цьому випадку проблему можна розв'язати методом розділення змінних (див. перший приклад з цієї лекції). Якщо перетворене рівняння виявиться гетерогенним, тоді нову задачу потрібно розв'язати іншим методом.

2. Лінійні неоднорідні граничні умови найзагальної форми
\end{verbatim}

\texttt{\textbackslash{}alpha\_1\ u\_x(0,\ t)\ +\ \textbackslash{}beta\_1\ u(0,\ t)\ =\ g\_1(t),}

\begin{verbatim}
~~
\alpha_2 u_x(L, t) + \beta_2 u(L, t) = g_2(t) Ви також можете перетворити умови на нульові граничні умови, використовуючи метод, описаний у другому прикладі. Звісно, нове рівняння, ймовірно, буде гетерогенним.
~~

- 3. Деякі методи розв'язку не містять вимог щодо унікальності граничних умов і, отже, не потребують попередніх перетворень. Пізніше, коли ми ознайомимося з перетворенням Лапласа, побачимо, що нульові граничні умови є необов'язковими (у випадку нульових граничних умов розв'язок дещо простіший, ніж в інших випадках).

- 4. Для граничних умов типу
\end{verbatim}

u(0, t) = g\_1(t),

\begin{verbatim}
u(L, t) = g_2(t) Метод, описаний у другому прикладі, призводить до наступної форми розв'язання задач:
\end{verbatim}

\texttt{u(x,\ t)\ =\ \textbackslash{}left\textbackslash{}\{\ g\_1(t)\ +\ \textbackslash{}frac\{x\}\{L\}\ \textbackslash{}left{[}\ g\_2(t)\ -\ g\_1(t)\ \textbackslash{}right{]}\ \textbackslash{}right\textbackslash{}\}\ +\ U(x,\ t).}

\begin{verbatim}
# ЗАВДАННЯ

1. Знайти розв'язок наступної змішаної задачі:
\end{verbatim}

\texttt{\textbackslash{}begin\{array\}\{lll\}\ (\textbackslash{}text\{УЧП\})\ \&\ u\_t\ =\ \textbackslash{}alpha\^{}2\ u\_\{xx\}\ +\ x\ \textbackslash{}cos\ t,\ \&\ 0\ \textless{}\ x\ \textless{}\ 1,\ \textbackslash{}\textbackslash{}\ (\textbackslash{}Gamma\ \textbackslash{}text\{У\})\ \&\ \textbackslash{}begin\{cases\}\ u\ \textbackslash{}\ (0,\ \textbackslash{}\ t)\ =\ 1,\ \textbackslash{}\textbackslash{}\ u\_x\ \textbackslash{}\ (1,\ \textbackslash{}\ t)\ +\ hu\ \textbackslash{}\ (1,\ \textbackslash{}\ t)\ =\ 1,\ \textbackslash{}end\{cases\}\ \&\ 0\ \textless{}\ t\ \textless{}\ \textbackslash{}infty,\ \textbackslash{}\textbackslash{}\ (\textbackslash{}text\{HY\})\ \&\ u\ \textbackslash{}\ (x,\ \textbackslash{}\ 0)\ =\ \textbackslash{}sin\ (\textbackslash{}pi\ x),\ \&\ 0\ \textbackslash{}leqslant\ x\ \textbackslash{}leqslant\ 1,\ \textbackslash{}end\{array\}}

\begin{verbatim}
попередньо перетворюючи його на задачу з нульовими граничними умовами. Чи відповідає рішення, яке ви отримуєте, вашим задумам щодо цієї проблеми?

Завдання

(PPP)
\end{verbatim}

\texttt{u\_t\ =\ u\_\{xx\},\ \textbackslash{}quad\ 0\ \textless{}\ x\ \textless{}\ 1,}

\begin{verbatim}
(GI)
\end{verbatim}

\texttt{\textbackslash{}begin\{cases\}\ u(0,\ t)\ =\ 0,\ \&\ 0\ \textless{}\ t\ \textless{}\ \textbackslash{}infty,\ \textbackslash{}\textbackslash{}\ u(1,\ t)\ =\ 1,\ \&\ 0\ \textless{}\ t\ \textless{}\ \textbackslash{}infty,\ \textbackslash{}end\{cases\}}

\begin{verbatim}
~~
(НУ) u(x, 0) = x^2, \quad 0 \le x \le 1, Перетворіть у задачу з нульовими граничними умовами і розв'яжіть її. Як виглядатиме рішення у різних
~~

Миті часу? Чи відповідає це вашим інтуїтивним ідеям? Що таке стаціонарне рішення? Як виглядає перехідне рішення?

3. Завдання
\end{verbatim}

\texttt{\textbackslash{}begin\{aligned\}}
\texttt{(\textbackslash{}text\{УЧП\})\textbackslash{}quad\ \&\ u\_t\ =\ u\_\{xx\},\ \&\&\ 0\ \textless{}\ x\ \textless{}\ 1,\textbackslash{}\ 0\ \textless{}\ t\ \textless{}\ \textbackslash{}infty,\ \textbackslash{}\textbackslash{}}
\texttt{(\textbackslash{}text\{ГУ\})\textbackslash{}quad\ \&\ u\_x(0,t)\ =\ 0,\textbackslash{}quad\ u\_x(1,t)\ +\ h\ u(1,t)\ =\ 1,\ \&\&\ 0\ \textless{}\ t\ \textless{}\ \textbackslash{}infty,\ \textbackslash{}\textbackslash{}}
\texttt{(\textbackslash{}text\{НУ\})\textbackslash{}quad\ \&\ u(x,0)\ =\ \textbackslash{}sin(\textbackslash{}pi\ x),\ \&\&\ 0\ \textbackslash{}leqslant\ x\ \textbackslash{}leqslant\ 1.}
\texttt{\textbackslash{}end\{aligned\}} \$\$

Перетворимо на задачу з нульовими граничними умовами. Чи буде нове
рівняння однорідним?

Лекція 7

\section{РОЗВ'ЯЗАННЯ СКЛАДНІШИХ ЗАДАЧ ШЛЯХОМ РОЗДІЛЕННЯ
ЗМІННИХ}\label{ux440ux43eux437ux432ux44fux437ux430ux43dux43dux44f-ux441ux43aux43bux430ux434ux43dux456ux448ux438ux445-ux437ux430ux434ux430ux447-ux448ux43bux44fux445ux43eux43c-ux440ux43eux437ux434ux456ux43bux435ux43dux43dux44f-ux437ux43cux456ux43dux43dux438ux445}

МЕТА ЛЕКЦІЇ: Показати, як більш складні проблеми теплопровідності можна
розв'язати методом розділення змінних. У цій лекції ми головним чином
розглядаємо такі задачі, розв'язання яких дозволить читачеві краще
ознайомитися з методом ділення змінних. Сподіваємося, що читач зможе
екстраполювати викладені тут ідеї для власних завдань.

У тій же лекції читач ознайомиться з задачею за її власними функціями,
або, як її ще називають, проблемою Штурма-Ліувіля. Обговорюються деякі
властивості проблеми Штурма-Ліувіля.

Як згадувалося вище, мета цієї лекції --- розв'язати змішану задачу
шляхом розділення змінних. Можливо, саме таким методом читач вирішить
свою проблему. Водночас ми сподіваємося, що він зможе перенести метод,
розглянутий у цій лекції, для власних завдань.

Почнемо з одномірної задачі теплопровідності, в якій одна з граничних
умов містить похідну.

\end{document}
